% 金星（综述）
% license CCBYSA3
% type Wiki

本文根据 CC-BY-SA 协议转载翻译自维基百科\href{https://en.wikipedia.org/wiki/Venus}{相关文章}。

\begin{figure}[ht]
\centering
\includegraphics[width=8cm]{./figures/044abc6191c52b10.png}
\caption{金星图像：由水手 10 号拍摄的两张照片合成，一张使用了清晰滤镜，另一张使用了蓝色滤镜，并经过平衡处理以在整体略带黄色的前提下逼近真实色彩\(^\text{[1]}\)。金星表面的云层是永久遮蔽的。} \label{fig_Venus_33}
\end{figure}
金星是距离太阳第二近的行星。由于其轨道最接近地球，且两者同为类地岩石行星，并在大小、质量和表面引力上几乎完全一致，因此金星常被称为地球的“孪生”或“姐妹”行星。然而，金星在许多方面与地球显著不同，尤其是它没有液态水，其大气层远比太阳系中其他任何岩石天体都更加厚重和致密。金星大气主要由二氧化碳组成，并覆盖着一层浓密的硫酸云层，环绕整个行星。在平均地表高度处，大气温度达到 737\,K（464\,°C；867\,°F），压力为地球海平面的 92 倍，使其最低层大气呈现超临界流体状态。从地球观测时，金星呈现为类似恒星的亮点，比天空中任何其他天然光源都更明亮\(^\text{[22][23]}\)，并且作为一颗内行星，它总是出现在接近太阳的方向，要么作为最明亮的“启明星”，要么作为“长庚星”。

金星与地球的轨道使两者在约 1.6 年的会合周期中彼此接近。在此过程中，金星会比其他任何行星更接近地球；相较之下，由于水星的轨道更靠近太阳，因此它在自身轨道路径中整体上比其他行星更接近地球。在从地球出发的行星际空间飞行中，金星常被用作引力助推的中转点，从而提供更快速且更经济的航程。金星没有卫星，并以非常缓慢的逆行方式绕其自转轴旋转，这被认为是太阳潮汐锁定效应与金星巨大大气层的受热差异相互竞争的结果。因此，金星的一昼夜长达 116.75 个地球日，而其一个太阳年为 224.7 个地球日。

金星具有微弱的磁层；由于缺乏内部发电机，其磁场是由太阳风与大气相互作用所诱导产生的。在内部结构上，金星具有核心、地幔和地壳。内部热量通过活跃的火山活动逸出\(^\text{[24][25]}\)，从而导致地表更新，而不是板块构造。金星在其早期历史中可能曾拥有液态地表水及宜居环境\(^\text{[26][27]}\)，但由于失控的温室效应，所有水分最终蒸发，使金星演变为当前的状态\(^\text{[28][29][30]}\)。在云层高度处的大气条件是整个太阳系中与地球最为相似的，这些条件被认为可能有利于金星生命的存在，并在 2020 年发现了潜在的生物标志物，从而推动了新的研究与探测任务。

在人类历史上，全世界的人们都曾观测过金星，并使其在众多文化中具有特殊的重要性。随着望远镜的出现，金星的相位变得可辨识，并在 1613 年被呈现为推翻当时占主导地位的地心模型、支持日心模型的决定性证据。1961 年，“金星 1 号”首次造访金星，飞掠行星并实现了人类第一次行星际飞行。1962 年，第二次行星际任务“水手 2 号”带回了首批来自金星的数据。1967 年，首个行星际撞击器“金星 4 号”抵达金星，随后在 1970 年由着陆器“金星 7 号”成功登陆。截至 2025 年，“太阳轨道飞行器”正前往于 2026 年飞掠金星，而下一项计划发射前往金星的任务是“金星生命探测器”，预定同样于 2026 年发射。
\subsection{物理特征}
\begin{figure}[ht]
\centering
\includegraphics[width=8cm]{./figures/23fb842f59355847.png}
\caption{金星（从左数第二，伪彩色显示）按比例与太阳系内侧具有行星质量的天体一起展示，排列顺序依其轨道自太阳向外依次排列（从左到右：水星、金星、地球、月球、火星和谷神星）。} \label{fig_Venus_1}
\end{figure}
\begin{figure}[ht]
\centering
\includegraphics[width=8cm]{./figures/f163eb810133acdf.png}
\caption{金星在不同波长下的成像} \label{fig_Venus_2}
\end{figure}
金星是太阳系中四颗类地行星之一，这意味着它与地球一样属于岩石天体。它在大小和质量上与地球相似，因此常被描述为地球的“姐妹”或“孪生”行星\(^\text{[31]}\)。由于其自转缓慢，金星的形状非常接近球形\(^\text{[32]}\)。其直径为 12,103.6\,km（7,520.8\,mi），仅比地球小 638.4\,km（396.7\,mi）；其质量为地球的 81.5\%，使金星成为太阳系中第三小的行星。金星表面的环境与地球完全不同，因为其致密大气中 96.5\% 为二氧化碳，导致强烈的温室效应，其余约 3.5\% 为氮气\(^\text{[33]}\)。金星表面气压为 9.3\,MPa（93\,bar），平均表面温度为 737\,K（464\,°C；867\,°F），高于两种主要组分的临界点，使其地表大气成为一种超临界流体，主要由超临界二氧化碳及部分超临界氮组成。
\subsubsection{自然历史}
\textbf{形成}

包括金星在内的类地岩石行星被认为经历了五个阶段形成：尘埃沉降、微行星形成、行星胚胎阶段、巨撞阶段，以及最终的大气形成阶段。来自金星的有限测量数据使其形成时间线难以进行更详细的分析\(^\text{[34]}\)。

\textbf{未来}

金星预计将在七至八十亿年后当太阳演变为红巨星时，与水星一起被毁灭，也有可能包括地球和月球\(^\text{[35]}\)。
\subsubsection{地理}
\begin{figure}[ht]
\centering
\includegraphics[width=8cm]{./figures/884bdec4a9f38b8d.png}
\caption{彩色编码的高程地图，显示金星上呈黄色的高地“大陆”以及次要地形特征} \label{fig_Venus_3}
\end{figure}
\begin{figure}[ht]
\centering
\includegraphics[width=8cm]{./figures/4ef3c5c0b6c896cf.png}
\caption{金星表面雷达数据的球形视图，突出显示地表特征（1989 年，麦哲伦号）。这些颜色并不代表地表的真实外观[36]。} \label{fig_Venus_4}
\end{figure}
在 20 世纪探测器揭示部分秘密之前，金星表面一直是推测的对象。1975 年和 1982 年的“金星号”着陆器传回了由沉积物和相对棱角分明的岩石覆盖的地表图像\(^\text{[37]}\)。1990–91 年间，“麦哲伦号”对其表面进行了详细测绘。有大量火山活动的证据，并且大气中二氧化硫的变化可能表明存在活火山\(^\text{[38][39]}\)。

大约 80\% 的金星表面被光滑的火山平原覆盖，其中 70\% 为带有皱褶脊的平原，10\% 为光滑或叶状平原\(^\text{[40]}\)。其余的地表由两个高地“大陆”构成，一个位于行星的北半球，另一个位于赤道以南。北部高地称为伊什塔尔高原（Ishtar Terra），得名于巴比伦的爱神伊什塔尔，其面积约与澳大利亚相当。麦克斯韦山脉位于伊什塔尔高原，其最高峰斯卡迪山（Skadi Mons）是金星上的最高点，海拔比金星平均表面高度高 11\,km（7\,mi）\(^\text{[41]}\)。南部高地称为阿芙洛狄忒高原（Aphrodite Terra），得名于希腊神话中的爱神，是两大高地中较大的一个，面积约与南美洲相当。该区域的大部分被裂缝和断层网络覆盖\(^\text{[42]}\)。

近期有关于金星上岩浆流动的证据（2024 年）\(^\text{[43]}\)，例如在盾状火山 Sif Mons 的熔岩流以及在平坦平原 Niobe Planitia 上的熔岩流\(^\text{[24]}\)。地表可见破火山口。金星的撞击坑数量很少，表明其表面相对年轻，约为 3 亿至 6 亿年\(^\text{[44][46]}\)。金星除了拥有岩石行星上常见的撞击坑、山脉和山谷外，还具有一些独特的地表特征。其中包括称为“farra”的平顶火山构造，其形状类似煎饼，直径从 20 至 50\,km（12 至 31\,mi）不等，高度在 100 至 1,000\,m（330 至 3,280\,ft）之间；称为“novae”的放射状、星形裂缝系统；具有放射状和同心裂隙、类似蜘蛛网的构造，被称为“arachnoids”；以及“coronae”，即被洼地包围的环状裂缝结构。这些特征均具有火山成因\(^\text{[46]}\)。
\begin{figure}[ht]
\centering
\includegraphics[width=8cm]{./figures/5421d90f38272eab.png}
\caption{彩色化影像（“金星 9 号”，1975 年），金星天空在地表呈橙黄色，这是由于瑞利散射或低层大气中的蓝色吸收体所致，而在更高的高度则呈白色[47][48]；同时其地表为类似玄武岩的深灰色，可能因氧化而呈红色[36]。} \label{fig_Venus_5}
\end{figure}
大多数金星表面特征以历史和神话中的女性命名\(^\text{[49]}\)。例外包括麦克斯韦山脉，以詹姆斯·克拉克·麦克斯韦命名，以及阿尔法高地区（Alpha Regio）、贝塔高地区（Beta Regio）和奥夫达高地区（Ovda Regio）。这后三个特征是在国际天文学联合会（负责行星命名的机构）采用现行系统之前命名的\(^\text{[50]}\)。

金星地形特征的经度以其本初子午线为基准表达。最初的本初子午线通过位于阿尔法高地区南部、椭圆形地形“夏娃”（Eve）中央的雷达亮点\(^\text{[51]}\)。在“金星号”任务完成后，本初子午线被重新定义为通过塞德娜平原（Sedna Planitia）上阿里阿德涅（Ariadne）陨石坑中央峰的位置\(^\text{[9][52]}\)。

地层学上最古老的镶嵌地形（tessera terrains）在“金星快车号”和“麦哲伦号”测量中显示出相对于周围玄武质平原更低的热辐射率，这表明其具有不同的、可能更富长英质（felsic）的矿物组合\(^\text{[29][53]}\)。生成大量长英质地壳的机制通常需要存在海洋和板块构造，这意味着早期金星上可能存在宜居条件，并在某一时期拥有大量水体\(^\text{[54]}\)。然而，镶嵌地形的具体性质仍不确定\(^\text{[55]}\)。

2023 年的研究首次提出金星在古代可能存在板块构造，因此可能拥有更宜居的环境，甚至可能具备维持生命的能力\(^\text{[26][27]}\)。金星因此成为研究类地行星形成及其宜居性的重要案例。

\textbf{火山活动}
\begin{figure}[ht]
\centering
\includegraphics[width=8cm]{./figures/96e9067c358877bf.png}
\caption{金星埃斯特拉区域内两个煎饼穹丘的雷达拼接图——二者均宽 65\,km（40\,mi），高度不足 1\,km（0.62\,mi）} \label{fig_Venus_6}
\end{figure}
金星表面的很大一部分似乎是由火山活动塑造的。金星上的火山数量是地球的数倍，并拥有 167 座直径超过 100\,km（60\,mi）的大型火山。地球上唯一与其规模相当的火山构造是夏威夷大岛\(^\text{[46]:154}\)。金星上已识别并绘制地图的火山数量超过 85,000 座\(^\text{[56][57]}\)。这并不是因为金星的火山活动比地球更活跃，而是因为其地壳更古老，且没有经历地球上活跃的侵蚀过程。地球的海洋地壳在构造板块边界通过俯冲不断被循环，其平均年龄约为 1 亿年\(^\text{[58]}\)，而金星表面估计为 3 亿至 6 亿年\(^\text{[44][46]}\)。

多项证据指向金星上持续存在火山活动。上层大气中的二氧化硫浓度在 1978 年至 1986 年间下降了 10 倍，在 2006 年急剧上升，然后再次下降 10 倍\(^\text{[59]}\)。这可能意味着多次大型火山喷发提高了其含量\(^\text{[60][61]}\)。有人提出金星的闪电（见下文）可能源自火山活动（即火山闪电）。2020 年 1 月，天文学家报告发现证据表明金星目前存在火山活动，具体而言，他们检测到橄榄石这种火山产物，而该物质在金星表面会迅速风化\(^\text{[62][63]}\)。

这一巨大的火山活动由炽热的内部提供能量，根据模型，它可以由行星早期阶段的高能碰撞以及类似地球的放射性衰变来解释。撞击的速度会明显高于地球，一方面因为金星更靠近太阳而运动更快，另一方面因为高偏心率天体与金星相撞时具有更高速度\(^\text{[64]}\)。

在 2008 年和 2009 年，“金星快车号”首次观测到正在进行火山活动的直接证据，其形式是在裂谷区加尼斯峡谷（Ganis Chasma）内的四个短暂出现的局部红外热点\(^\text{[note 1]}\)，该区域靠近盾状火山玛特山（Maat Mons）\(^\text{[65]}\)。其中三个热点在连续轨道中被重复观测到。这些热点被认为代表新近喷发的熔岩\(^\text{[66][67]}\)。实际温度未知，因为热点的面积无法测量，但推测其温度可能在 800–1,100\,K（527–827\,°C；980–1,520\,°F）范围内，而正常温度约为 740\,K（467\,°C；872\,°F）\(^\text{[68]}\)。2023 年，科学家重新分析了“麦哲伦号”拍摄的玛特山区域地形图像。通过计算机模拟，他们确认该区域在 8 个月间发生了地形变化，并得出其原因是活跃火山活动\(^\text{[69]}\)。

\textbf{陨石坑}
\begin{figure}[ht]
\centering
\includegraphics[width=8cm]{./figures/d8116d164677ab95.png}
\caption{金星表面的撞击坑（伪彩色，由雷达数据重建的三维投影影像）} \label{fig_Venus_7}
\end{figure}
金星上有将近一千个撞击坑，均匀分布在其表面。在其他具有撞击坑的天体上，例如地球和月球，撞击坑呈现出不同程度的退化状态。月球上的退化由后续撞击造成，而地球上的退化则由风化和降雨侵蚀导致。在金星上，约 85\% 的撞击坑保持完好状态。撞击坑的数量及其良好的保存状况表明，该行星曾在 3 亿至 6 亿年前经历过一次全球性地表重塑事件\(^\text{[44][45]}\)，随后火山活动逐渐衰减\(^\text{[70]}\)。地球的地壳处于持续运动中，而金星被认为无法维持此类过程。在没有板块构造来散发地幔热量的情况下，金星经历一种循环过程，即地幔温度不断升高，直到达到削弱地壳的临界水平。随后，在大约 1 亿年的时间内，发生大规模的俯冲，将地壳完全循环\(^\text{[46]}\)。

金星的撞击坑直径范围从 3 到 280\,km（2 到 174\,mi）。没有直径小于 3\,km 的撞击坑，这是由于致密大气对进入物体的影响所致。动能低于某一阈值的物体会被大气大幅减速，以至于无法形成撞击坑\(^\text{[71]}\)。直径小于 50\,m（160\,ft）的入射物体将在到达地面之前于大气层中解体并焚毁\(^\text{[72]}\)。
\subsubsection{内部结构}
\begin{figure}[ht]
\centering
\includegraphics[width=8cm]{./figures/dbfa21d636bc841c.png}
\caption{金星的分化结构} \label{fig_Venus_8}
\end{figure}
由于缺乏反射地震学的数据或其转动惯量的测量结果，人们对金星内部结构和地球化学性质的直接信息非常有限\(^\text{[73]}\)。金星与地球在大小和密度上的相似性表明，它们可能拥有相似的内部结构：核心、地幔和地壳。与地球相似，金星的核心很可能至少部分为液态，因为这两颗行星的冷却速率大致相同\(^\text{[74]}\)，尽管完全固态核心也不能被排除\(^\text{[75]}\)。金星体积略小意味着其深部内部的压强比地球低 24\%\(^\text{[76]}\)。根据行星模型预测的转动惯量数值，金星核心半径估计为 2,900–3,450\,km\(^\text{[75]}\)。基于 2006 至 2020 年之间测量的自转轴进动速率所推算的转动惯量，目前的估计为 3,500\,km\(^\text{[13][77]}\)。

金星地壳平均估计厚度为 40\,km，最大可达 65\,km\(^\text{[78]}\)。

两颗行星的主要差异在于金星缺乏板块构造的证据，可能原因是其地壳太坚固，在没有水减少黏滞性的情况下无法发生俯冲。这导致行星热量损失较少，阻碍其冷却，并为其缺乏内部磁场提供了可能的解释\(^\text{[79]}\)。金星可能改以周期性的全球地表重塑事件释放内部热量\(^\text{[44]}\)。
\subsubsection{磁场与核心}
\begin{figure}[ht]
\centering
\includegraphics[width=8cm]{./figures/81a077950d438b05.png}
\caption{} \label{fig_Venus_9}
\end{figure}
1967 年，“金星 4 号”发现金星的磁场比地球弱得多。该磁场是由电离层与太阳风之间的相互作用所诱导产生的\(^\text{[80][81][page needed]}\)，而不是像地球那样由核心内部的发电机产生。金星这种微弱的诱导磁层对大气抵御太阳和宇宙辐射的保护几乎可以忽略不计。

金星缺乏本征磁场令科学家感到意外，因为金星在大小上与地球相似，原本应在其核心中存在发电机。发电机需要三个条件：导电液体、自转和对流。金星的核心被认为具有电导性，而尽管其自转通常被认为过于缓慢，模拟表明其仍足以产生发电机\(^\text{[82][83]}\)。这意味着缺乏发电机的原因在于金星核心内部缺乏对流。在地球上，对流发生在核心的液态外层中，因为液态层底部的温度远高于顶部。在金星上，一次全球性的地表重塑事件可能终止了板块构造，导致通过地壳的热通量减小。这种隔热效应会使地幔温度升高，从而减少核心的热通量。因此，无法提供用于驱动磁场的内部地球发电机，核心释放的热量则重新加热地壳\(^\text{[84]}\)。

一种可能性是金星没有固态内核\(^\text{[85]}\)，或者其核心没有在冷却，使得整个液态核心的温度大致相同。另一种可能性是其核心已经完全固化。核心的状态高度依赖于硫的浓度，而当前尚不清楚\(^\text{[84]}\)。

另一种可能性是金星缺乏像地球“形成月球”事件那样的大型撞击，使金星核心在增量形成过程中产生分层结构，而没有足够的力量触发或维持对流，从而无法形成“地质发电机”\(^\text{[86]}\)。

金星周围微弱的磁层意味着太阳风会直接与其外层大气相互作用。在这里，水分子因紫外线辐解而生成氢和氧离子。随后，太阳风提供能量，使其中一些离子获得足够速度逃离金星的引力场。该侵蚀过程导致轻质的氢、氦和氧离子持续流失，而高质量分子如二氧化碳更可能被保留。太阳风导致的大气侵蚀可能造成金星在形成后的前十亿年内失去了大部分水\(^\text{[87]}\)。然而，该行星在最初的 20–30 亿年内可能保有发电机，因此水的流失也可能发生得更晚\(^\text{[88]}\)。这种侵蚀使得大气中高质量氘与低质量氢的比值相比太阳系其他区域增加了 100 倍\(^\text{[89]}\)。
\subsection{大气与气候}
\begin{figure}[ht]
\centering
\includegraphics[width=8cm]{./figures/8e5406aad8f1ba0c.png}
\caption{通过紫外成像显现的金星大气云层结构} \label{fig_Venus_10}
\end{figure}
金星拥有致密的大气层，由 96.5\% 的二氧化碳、3.5\% 的氮气组成——两者在金星地表均以超临界流体形式存在，其密度为水的 6.5\%\(^\text{[90]}\)——并含有微量其他气体，包括二氧化硫\(^\text{[91]}\)。其大气质量为地球的 92 倍，而地表压力约为地球的 93 倍，相当于地球海洋表面下近 1\,km（5⁄8\,mi）处的压力。地表大气密度为 65\,kg/m³（4.1\,lb/cu\,ft），为水的 6.5\%\(^\text{[90]}\)，或为地球海平面 293\,K（20\,°C；68\,°F）空气密度的 50 倍。富含 CO₂ 的大气产生了太阳系中最强的温室效应，使其表面温度至少达到 735\,K（462\,°C；864\,°F）\(^\text{[92][93]}\)。这使得金星表面比水星更热\(^\text{[94]}\)，尽管金星与太阳的距离几乎是水星的两倍，因此其接收的太阳辐照度仅为水星的四分之一，即 2,600\,W/m²（约为地球的两倍）\(^\text{[5]}\)。由于其失控温室效应，金星被卡尔·萨根等科学家视为与地球气候变化相关的警示和研究对象\(^\text{[95]}\)。因此金星被称为温室行星\(^\text{[96]}\)，或处于温室炼狱中的行星\(^\text{[97]}\)。

金星大气中富含原初稀有气体，与地球大气相比更为丰富\(^\text{[98]}\)。这种富集表明金星在演化早期便与地球发生分歧。一个异常大的彗星撞击\(^\text{[99]}\)或从太阳星云吸积更大质量的原始大气\(^\text{[100]}\)被提出用于解释这种富集。然而，大气中放射性氩–40 含量较低，这种元素可视为地幔脱气的替代指标，这表明主要岩浆作用在早期便已终止\(^\text{[101][102]}\)。

研究表明，在数十亿年前，金星的大气可能与早期地球的大气非常相似，并且其地表可能存在大量液态水\(^\text{[103][104][105]}\)。经过约 6 亿年至数十亿年的时间\(^\text{[106]}\)，太阳光度上升及可能发生的大规模火山重塑导致原始水体蒸发\(^\text{[107]}\)。当大气中的温室气体（包括水）达到临界值后，便触发了失控温室效应\(^\text{[108]}\)。尽管金星地表环境已不再适合任何可能在此事件前形成的类地生命，但人们推测金星表面上方 50\,km（30\,mi）的高层云层中可能存在生命，因为该高度的环境是整个太阳系中最类似地球的\(^\text{[109]}\)，温度介于 303 至 353\,K（30 至 80\,°C；86 至 176\,°F）之间，气压和辐射水平与地球表面相当，但大气成分为酸性云和二氧化碳\(^\text{[110][111][112]}\)。更具体而言，在 48 至 59\,km 的高度范围内，温度和辐射条件适宜生命生存；在更低的高度水会蒸发，而在更高的高度紫外辐射将过强\(^\text{[113][114]}\)。2020 年 9 月，有研究报告在金星大气中探测到磷化氢的吸收线，而该物质在已知条件下无法通过非生物途径生成，因此引发了关于金星大气中可能存在现存生命的猜测\(^\text{[115][116]}\)。后续研究将这一光谱信号解释为二氧化硫\(^\text{[117]}\)，或发现实际上不存在该吸收线\(^\text{[118][119]}\)。
\begin{figure}[ht]
\centering
\includegraphics[width=8cm]{./figures/b77f963239de7576.png}
\caption{按高度绘制的大气剖面图（左侧刻度）：云层、温度变化（VIRA 粗线及底部刻度）、压力变化（右侧刻度）以及风速（PV 虚线及顶部刻度）} \label{fig_Venus_11}
\end{figure}
热惯性以及下层大气中由风输送热量的作用意味着，尽管金星自转缓慢，其面对太阳与背向太阳的半球之间的地表温度并没有显著差异。地表的风速很慢，仅为每小时数公里，但由于地表大气密度极高，它们对障碍物施加的力非常大，并能够在地表搬运尘埃和小石块。即便不考虑其高温、高压和缺氧环境，这种阻力本身就足以使人在地表难以行走\(^\text{[120]}\)。

在致密的 CO₂ 层之上，距地表 45 至 70\,km 的高度存在厚厚的云层\(^\text{[121]}\)，主要由硫酸组成，其形成过程是二氧化硫分子在紫外辐射催化下与水发生反应\(^\text{[122]}\)，从而生成硫酸水合物\(^\text{[123]}\)。此外，云层中含有约 1\% 的三氯化铁\(^\text{[124][125]}\)。云粒子的其他可能成分包括硫酸铁、氯化铝和五氧化二磷。不同高度的云层具有不同的成分和粒径分布\(^\text{[124]}\)。这些云层像地球的厚云层一样，将落在其上的阳光约 70\% 反射回太空\(^\text{[1265]}\)，并由于覆盖整个行星，使得金星表面无法被可见光直接观测。永久性的云层意味着尽管金星比地球更接近太阳，其地表接收到的阳光却更少，只有 10\% 的入射阳光能到达地表\(^\text{[127]}\)，使得其地表白天平均照度约为 14,000 勒克斯，与地球“阴天白昼”的照度相当\(^\text{[128]}\)。

金星大气的自转速度远快于其固体本体的自转，这一现象被称为大气超自转\(^\text{[129]}\)。这导致云顶处出现强劲的 300\,km/h（185\,mph）风，它们绕行星一周仅需约 4 天，相当于行星自转速度的 60 倍\(^\text{[130]}\)，而地球上最强的风速仅为其自转速度的 10–20\%。

尽管金星在可见光下看起来缺乏表面特征，但在紫外波段却存在条带或条纹，其成因尚未确定。紫外吸收可能来源于一种由氧与硫组成的化合物 OSSO，该分子中两个硫原子之间具有双键，并存在“顺式”和“反式”两种构型，或可能来源于 S₂ 至 S₈ 的多硫化合物\(^\text{[131]}\)。
\begin{figure}[ht]
\centering
\includegraphics[width=8cm]{./figures/d0e3e9800a1218c4.png}
\caption{} \label{fig_Venus_12}
\end{figure}
金星的地表实际上是等温的；它不仅在两个半球之间保持恒定温度，而且在赤道与两极之间也保持恒定温度\(^\text{[5][131]}\)。金星的自转轴倾角极小——不到 3°，相比之下地球为 23°——这也使季节性温度变化极小\(^\text{[134]}\)。海拔是影响金星温度为数不多的因素之一。金星最高点——麦克斯韦山脉中的斯卡迪山——因此是金星上最冷的地点，温度约为 655\,K（380\,°C；715\,°F），气压约为 4.5\,MPa（45\,bar）\(^\text{[132][135]}\)。1995 年，“麦哲伦号”探测器在金星最高山峰顶部成像到一种高反射物质，被称为“金星雪”，与地球上的雪极为相似。这种物质可能通过类似于雪的过程形成，但温度高得多。由于其挥发性过高而无法在地表凝结，它以气态上升至更高、更冷的高度，在那里可能发生沉降。该物质的成分尚未确定，但推测范围从碲单质到硫化铅（方铅矿）不等\(^\text{[136]}\)。

尽管金星没有季节变化，天文学家在 2019 年发现其大气对阳光吸收存在周期性变化，可能由悬浮在高层云中的不透明吸收粒子引起。这种变化导致金星纬向风速出现观测到的变化，并似乎与太阳 11 年黑子周期同步\(^\text{[137]}\)。

金星大气中是否存在闪电长期以来存在争议\(^\text{[138]}\)，自从苏联“金星号”探测器首次探测到疑似闪电脉冲后\(^\text{[139][140][141]}\)。在 2006–07 年，“金星快车号”清晰探测到了哨声模式波，这是闪电的特征信号。其间歇性出现表明其与天气活动相关。根据这些测量，其闪电发生率至少为地球的一半\(^\text{[142]}\)，然而其他仪器完全未能探测到闪电\(^\text{[139]}\)。闪电的起源仍不明确，但可能来自云层或金星火山。

2007 年，“金星快车号”发现金星南极存在一个巨大的双极大气涡旋\(^\text{[143][144]}\)。“金星快车号”在 2011 年发现金星高层大气中存在臭氧层\(^\text{[145]}\)。2013 年，欧洲航天局科学家报告称金星的电离层向外流动，其方式类似于“在类似条件下彗星电离尾的外泄”\(^\text{[146][147]}\)。

在 2015 年 12 月，以及较弱的 2016 年 4 月和 5 月，参与日本“晓号”任务的研究人员在金星大气中观测到弧形结构。这被认为是太阳系中最大、可能也是首个直接证据确认的驻波重力波\(^\text{[148][149][150]}\)。

金星地表大气的颜色与声音\(^\text{[151]}\)已被记录：天空在地表呈橙黄色，而在更高的高度则呈白色\(^\text{[48]}\)。
\subsection{轨道与自转}
金星以约 0.72\,AU（1.08 亿\,km；6700 万\,mi） 的平均距离绕太阳运行，每 224.7 天完成一周。它在 7.998 年内完成 13 圈公转，因此在天空中的位置几乎每八年重复一次。尽管所有行星轨道都是椭圆形的，但金星的轨道目前最接近圆形，其偏心率小于 0.01\(^\text{[5]}\)。对早期太阳系轨道动力学的模拟显示，金星轨道的偏心率在过去可能大得多，高达 0.31，并可能影响其早期气候演化\(^\text{[152]}\)。
\begin{figure}[ht]
\centering
\includegraphics[width=8cm]{./figures/cefd4e3932d4f154.png}
\caption{金星及其自转相对于公转的关系} \label{fig_Venus_13}
\end{figure}
金星具有逆行自转，这意味着它绕自身轴旋转的方向与包括地球在内的大多数行星不同，为顺时针方向，与其绕太阳的逆时针公转方向相反。因此金星的恒星日为 243 个地球日，比其公转周期 224.7 个地球日更长。如果金星被潮汐锁定在太阳上，它将始终以同一面朝向太阳，其恒星日将为 224.7 天。然而，由于金星大气质量巨大且接近太阳，大气的差异加热使金星具有微弱的逆行自转。由于同样的原因，其昼长也可波动长达 20 分钟\(^\text{[153][154]}\)。利用“麦哲伦号”探测器数据在 500 天期间测量得到的金星自转周期小于在“麦哲伦号”和“金星快车号”访期间隔 16 年中测得的自转周期，两者相差约 6.5 分钟\(^\text{[155]}\)。由于逆行自转，金星的太阳日显著短于其恒星日，仅为 116.75 个地球日\(^\text{[12]}\)。一个金星年约相当于 1.92 个金星太阳日\(^\text{[156]}\)。对于位于金星地表的观察者而言，太阳将从西方升起并在东方落下\(^\text{[156]}\)，尽管金星不透明的云层使地表无法看到太阳\(^\text{[157]}\)。

金星可能在太阳星云中形成时具有不同的自转周期和自转轴倾角，并由于行星摄动和潮汐效应对其致密大气产生的混沌自转变化而达到当前状态，这一演变可能历时数十亿年。金星的自转周期可能代表一种平衡状态，即太阳引力的潮汐锁定效应（倾向于减慢自转）与厚重大气因太阳加热产生的大气潮汐之间的竞争\(^\text{[158][159]}\)。金星与地球连续两次接近之间的平均间隔为 584 天，这几乎正好等于 5 个金星太阳日（精确值为 5.001444）\(^\text{[160]}\)，但与地球存在自转-公转共振的假说已被否定\(^\text{[161]}\)。

金星没有天然卫星\(^\text{[162]}\)。它有数颗特洛伊小行星：拟卫星 524522 Zoozve\(^\text{[163][164]}\) 以及另外两颗临时特洛伊小行星 2001\,CK32 和 2012\,XE133\(^\text{[165]}\)。17 世纪，乔瓦尼·卡西尼报告称有一颗月亮绕金星运行，并将其命名为 Neith，接下来 200 年中出现了多次观测报告，但大多数后来被确定为附近恒星。加州理工学院 Alex Alemi 和 David Stevenson 在 2006 年对早期太阳系模型的研究显示，金星在数十亿年前一次巨大撞击事件中很可能形成过至少一颗卫星\(^\text{[166]}\)。根据该研究，大约 1000 万年后，另一场撞击事件使金星自转方向逆转，随后的潮汐减速使金星的卫星逐渐向内螺旋并最终与金星相撞\(^\text{[167]}\)。如果后续撞击再次形成过卫星，它们会以同样方式消失。另一种解释金星缺乏卫星的可能性是强烈的太阳潮汐作用会使内侧类地行星周围的大型卫星轨道失稳\(^\text{[162]}\)。

金星轨道空间中存在一条尘埃环云\(^\text{[168]}\)，其可能的来源包括金星追踪的小行星\(^\text{[169]}\)、呈波状迁移的行星际尘埃，或太阳系形成时原始环境盘的遗留物\(^\text{[170]}\)。
\subsubsection{相对于地球的轨道}
\begin{figure}[ht]
\centering
\includegraphics[width=8cm]{./figures/8445fe95a61e302d.png}
\caption{地球位于图示中央，曲线表示金星随时间变化的方向和距离。} \label{fig_Venus_14}
\end{figure}
地球与金星具有近似的 13:8 轨道共振（地球每公转八次，金星公转十三次）\(^\text{[171]}\)。因此，它们彼此接近并在平均 584 天的会合周期中达到下合\(^\text{[5]}\)。从地心视角看，金星相对于地球的轨迹在五个会合周期中绘制出一个五角星形，每个周期旋转 144°。由于该轨迹在视觉上类似花朵，这个金星五角星有时被称为金星之花瓣\(^\text{[172]}\)。

当金星在下合位置位于地球与太阳之间时，它在所有行星中以平均 4100 万\,km（2500 万\,mi）的距离最接近地球\(^\text{[5][note 2][173]}\)。由于地球轨道偏心率正在减小，这一最小距离将在未来数万年中逐渐增大。从公元 1 年到 5383 年，共出现 526 次距离小于 4000 万\,km（2500 万\,mi）的接近；之后约 60,158 年内将不再出现此类接近\(^\text{[174]}\)。

虽然金星接近地球的最近距离最小，但水星更频繁地成为距离地球（以及任何其他行星）最近的行星\(^\text{[175][176]}\)。金星已被用于引力助推的中转点，被认为是前往水星\(^\text{[177][178]}\)、太阳\(^\text{[179]}\)、小行星\(^\text{[180]}\)、火星\(^\text{[181]}\)、木星及更远处\(^\text{[182][183]}\)更快速且经济的航程方式。

在潮汐作用方面，金星对地球施加的潮汐力在月球与太阳之后位列第三，但要弱得多\(^\text{[184]}\)。
\subsection{可观测性}
\begin{figure}[ht]
\centering
\includegraphics[width=8cm]{./figures/149b4a28cac2124e.png}
\caption{金星（图中偏右的位置）从地球观测时在其最大亮度下始终比所有其他行星或恒星更亮。图像顶部可见木星。} \label{fig_Venus_15}
\end{figure}
在肉眼看来，金星呈现为一个白色光点，其最大视星等为 −4.92，比除太阳外的任何行星或恒星都要明亮\(^\text{[185]}\)，即使在凌日过程中最暗时，其视星等也为 −2.98\(^\text{[20]}\)。该行星的平均视星等为 −4.14，标准差为 0.31\(^\text{[20]}\)。在下合前后约一个月的娥眉月相期间，其亮度达到最大。当前方被太阳背光照亮时，金星会变暗至约 −3 等，但确切亮度取决于位相角\(^\text{[186]}\)。金星亮度足以在白昼中被看见\(^\text{[187]}\)，但在太阳位于地平线附近或正在落下时更易观测。作为一颗内行星，它与太阳的角距离始终不超过约 47°\(^\text{[188]}\)。

金星每 584 天在绕太阳运行时“追上”地球\(^\text{[5]}\)。在此过程中，它从日落后可见的“昏星”变为日出前可见的“晨星”。虽然另一颗内行星水星最大延伸角仅为 28°，在曙暮光中往往难以辨认，但金星在最亮时几乎不可能错过。其更大的最大延伸角意味着它在日落后仍能在暗空中被看到。作为夜空中最亮的点状天体，金星经常被误报为“不明飞行物”\(^\text{[189]}\)。

由于金星在下合时会靠近地球，且其轨道相对地球轨道有倾角，它可以出现在比黄道平面北或南超过 8° 的位置，这超过任何其他行星或月球。每八年左右的三月，金星会出现在距离黄道最北的位置，位于双鱼座（例如 2025 年 3 月中旬）；而每八年左右的八月或九月，它会出现在距离黄道最南的位置，位于室女座（例如 2023 年 8 月下旬）。因此在北半球，金星可能位于太阳以北，并在同一天既作为晨星又作为昏星出现。这些南北偏离的时间每年会略微提前，在 30 个周期（240 年）后，这一周期会逐渐被另一个相差三年的周期取代，因此在地球 243 次公转和金星 395 次公转后情况将接近恢复到原来的样子\(^\text{[190]}\)。

金星的月掩现象，即月球遮挡金星，使地球某些地区的观测者无法看到金星，平均每年大约发生两次，有时同一年会发生多次（尽管十分罕见）。
\subsubsection{相位}
\begin{figure}[ht]
\centering
\includegraphics[width=8cm]{./figures/7f0dd42656d60e11.png}
\caption{金星的相位及其视直径的变化} \label{fig_Venus_16}
\end{figure}
在绕太阳运行的过程中，金星在望远镜视野中呈现出类似月相的相位变化。当金星位于太阳的远端（上合）时，呈现为一个较小而“满”的圆面。在它距离太阳达到最大角距时，金星显示为更大的圆面并呈“上弦相”，此时它在夜空中最为明亮。当金星经过地球与太阳之间的近侧时，在望远镜中呈现为更大的细长“娥眉月”。当金星位于地球与太阳之间（下合）时，其视直径最大并呈“朔相”。其大气可通过望远镜观察到，这是由于阳光在其周围发生折射形成的光晕\(^\text{[188]}\)。这些相位在 4 英寸口径的望远镜中即可清晰可见\(^\text{[191]}\)。虽然肉眼能否观察到金星的相位存在争议，但仍有记录显示有人观察到其娥眉形状\(^\text{[192]}\)。
\subsubsection{日间显现}
\begin{figure}[ht]
\centering
\includegraphics[width=8cm]{./figures/b29b2a44adbea1a0.png}
\caption{金星在白昼中常可被肉眼看见，例如在 2015 年 12 月 7 日月掩事件发生前不久所观测到的情形。} \label{fig_Venus_17}
\end{figure}
当金星足够明亮且与太阳具有足够大的角距离时，它在晴朗的白昼天空中可被肉眼轻易看到，尽管大多数人并不知道要去寻找它\(^\text{[193]}\)。天文学家爱德蒙·哈雷在 1716 年计算了其肉眼可见的最大亮度，当时许多伦敦人因在白天看到金星而感到惊慌。法国皇帝拿破仑·波拿巴曾在卢森堡参加一次招待会时目睹过金星的白昼显现\(^\text{[194]}\)。另一次历史性的金星白昼观测发生在 1865 年 3 月 4 日美国总统亚伯拉罕·林肯在华盛顿特区举行就职典礼期间\(^\text{[195]}\)。
\subsubsection{凌日}
\begin{figure}[ht]
\centering
\includegraphics[width=6cm]{./figures/80b71a3a59e0792d.png}
\caption{通过望远镜投射到白色卡片上的 2012 年金星凌日影像} \label{fig_Venus_18}
\end{figure}
金星凌日是指在下合期间，金星出现在太阳前方。由于金星轨道相对于地球轨道略有倾角，因此大多数每隔 1.6 年发生的金星与地球的下合并不会产生金星凌日。因此，只有当下合发生在六月或十二月的若干天内，即金星和地球的轨道与太阳成一直线时，才会出现金星凌日\(^\text{[196]}\)。其结果是金星凌日目前呈现为一个 8 年、105.5 年、8 年和 121.5 年的序列，构成一个 243 年的周期。

金星凌日在历史上具有重要意义，因为它使天文学家能够确定天文单位的大小，从而推算出太阳系的规模。1639 年，杰里迈亚·霍罗克斯通过首次已知的金星凌日观测（继 1631 年人类首次观测到的行星凌日——水星凌日之后）展示了这一点\(^\text{[197]}\)。

自约翰内斯·开普勒在 1621 年计算出其发生时间以来，至今仅有七次金星凌日被观测到。库克船长于 1768 年航行至塔希提岛以记录第三次金星凌日，随后促成了对澳大利亚东海岸的探索\(^\text{[198][199]}\)。

最近的一对金星凌日发生在 2004 年 6 月 8 日和 2012 年 6 月 5–6 日。该凌日可通过许多在线平台进行实时观看，或在适当设备与条件下进行本地观测\(^\text{[200]}\)。再前一对则发生在 1874 年 12 月和 1882 年 12 月。

下一次金星凌日将发生在 2117 年 12 月和 2125 年 12 月\(^\text{[201]}\)。
\subsubsection{灰光现象}
\begin{figure}[ht]
\centering
\includegraphics[width=8cm]{./figures/28c0a4b35d424546.png}
\caption{自 2022 年起，夜辉被认为是灰光现象最可能的成因。在这张可见光与近红外影像中，它最容易被辨认为沿金星边缘的一条亮线\(^\text{[202]}\)。地表及其特征（如图中可见的暗色区域——阿芙洛狄忒高原的奥夫达高地区）在人眼中难以辨认得如此清晰，尽管据报道某些人能够看到，可能是因为他们对地表发光所在光谱具有更高的敏感度\(^\text{[203]}\)。} \label{fig_Venus_19}
\end{figure}
金星观测中的一个长期谜团是所谓的灰光——即在金星呈娥眉月相时，其暗面上出现的微弱亮光。首次关于灰光的观测声称可追溯至 1643 年，但这一亮光的存在从未得到可靠确认。观测者推测，这可能源于金星大气中的电活动，但也可能是错觉，由观察一个明亮的娥眉形天体所引发的生理效应造成\(^\text{[204][140]}\)。灰光现象常在金星出现在傍晚天空时被报告，此时金星的傍晚界线朝向地球。
\subsection{观测与探索历史}
\subsubsection{早期观测}
金星在地球天空中亮度足以裸眼可见，使其成为人类历史上各个文化所认识和识别的经典行星之一，尤其因为它是仅次于太阳和月亮的地球天空中第三亮的天体。由于金星的运动呈现不连续性（因其接近太阳而消失数日，然后在另一侧地平线上重新出现），一些文化没有将金星视为单一天体\(^\text{[205]}\)；相反，他们认为金星在东西地平线上分别是两颗不同的恒星：晨星与昏星\(^\text{[205]}\)。然而，杰姆代特·纳斯尔时期的圆筒印章以及巴比伦第一王朝的阿弥萨杜卡金星泥板都表明，古代苏美尔人已经知道晨星与昏星是同一天体\(^\text{[206][205][207]}\)。
\begin{figure}[ht]
\centering
\includegraphics[width=6cm]{./figures/07505049521d37a0.png}
\caption{已知最古老的金星位置记录，来自阿弥萨杜卡时期的巴比伦金星泥板（公元前 1600 年）} \label{fig_Venus_20}
\end{figure}
在古巴比伦时期，金星被称为 Ninsi'anna，后来被称为 Dilbat\(^\text{[208]}\)。“Ninsi'anna”一名意为“神圣的女士、天空的光辉”，指的是作为最亮可见“星辰”的金星。该名称更早的拼写使用楔形文字 si4（＝SU，意为“呈红色”），其原始含义可能是“天空之红的神圣女士”，指代晨空与昏空的颜色\(^\text{[209]}\)。

中国古代将晨星金星称为“太白”或“启明”，而将昏星金星称为“长庚”\(^\text{[210]}\)。

古希腊人最初认为金星是两颗独立的星：晨星 Phosphorus 与昏星 Hesperus。老普林尼将认识到这两者是同一天体的功劳归于公元前六世纪的毕达哥拉斯\(^\text{[211]}\)，而第欧根尼·拉尔修则认为应归功于巴门尼德（公元前五世纪早期）\(^\text{[212]}\)。尽管古希腊人已认识到金星是同一颗天体，古罗马人仍沿用不同称呼，将其晨星状态称为 Lucifer（字面意为“光之携带者”），将其昏星状态称为 Vesper\(^\text{[213]}\)，两者均为其传统希腊名称的直译。

在二世纪的《天文学大成》中，托勒密提出水星和金星均位于太阳与地球之间。十一世纪波斯天文学家阿维森纳声称曾观测到金星凌日（尽管这一点存在争议）\(^\text{[214]}\)，后来的天文学家将其视为托勒密理论的佐证\(^\text{[215]}\)。十二世纪安达卢西亚天文学家伊本·巴贾赫曾观测到“两个在太阳表面上的黑点”；十三世纪马拉盖天文学家库特布丁·希拉齐认为这些是金星和水星的凌日\(^\text{[216]}\)，但这不可能为真，因为在伊本·巴贾赫的一生中并无金星凌日\(^\text{[note 3]}\)。
\begin{figure}[ht]
\centering
\includegraphics[width=6cm]{./figures/1899ac8f6b035776.png}
\caption{前哥伦布时期的玛雅《德累斯顿抄本》，其中记录并计算了金星的出现时机} \label{fig_Venus_21}
\end{figure}
\subsubsection{金星与近代早期天文学}
\begin{figure}[ht]
\centering
\includegraphics[width=10cm]{./figures/258414620ee3be6f.png}
\caption{1610 年，伽利略通过望远镜观测到金星呈现出不同的相位，尽管它在地球天空中始终靠近太阳（第一幅图）。这一发现证明金星绕太阳运行而非绕地球运行，如哥白尼的日心模型所预言，并推翻了托勒密的地心模型（第二幅图）。} \label{fig_Venus_22}
\end{figure}
当金星于 1610 年 12 月首次被意大利物理学家伽利略用望远镜观测时，他发现金星呈现如月亮般的相位变化，从娥眉月到凸月再到满相，反之亦然。当金星在天空中远离太阳的位置时，它呈现半月相；而当它靠近太阳时，则呈现娥眉相或满相。只有当金星绕太阳运行时，这种现象才可能出现。伽利略在 1613 年的《太阳黑子书信》中报告了这一发现，这成为最早明确反驳托勒密地心体系（即太阳系同心且以地球为中心）的观测之一\(^\text{[219][220]}\)。

1631 年的金星凌日虽然未被记录下来，但这是首次被成功预测的凌日，由约翰内斯·开普勒计算并于 1629 年发表。接下来的 1639 年金星凌日由杰里迈亚·霍罗克斯精确预测，并由他与朋友威廉·克拉布特里各自在家中于 1639 年 12 月 4 日（当时使用的儒略历为 11 月 24 日）同时观测到\(^\text{[221]}\)。
\begin{figure}[ht]
\centering
\includegraphics[width=8cm]{./figures/802be5669f872ee3.png}
\caption{描绘杰里迈亚·霍罗克斯观测 1639 年金星凌日的二十世纪绘画} \label{fig_Venus_23}
\end{figure}
金星的大气由俄罗斯博学家米哈伊尔·罗蒙诺索夫于 1761 年发现\(^\text{[222][223]}\)。1790 年，德国天文学家约翰·施雷特尔观测了金星的大气。施雷特尔发现，当金星呈现细娥眉相时，其两端的尖角延伸超过 180°。他正确地推断这源于致密大气中阳光的散射。之后，美国天文学家切斯特·史密斯·莱曼在金星处于下合时观察到其暗面周围的完整光环，为其具有大气提供了进一步证据\(^\text{[224]}\)。大气的存在使确定金星自转周期的工作变得复杂，意大利出生的天文学家乔瓦尼·卡西尼和施雷特尔等观测者根据金星表观表面标记的运动错误地估计其自转周期约为 24 小时\(^\text{[225]}\)。
\begin{figure}[ht]
\centering
\includegraphics[width=6cm]{./figures/d02c63b404bf6fdf.png}
\caption{1769 年金星凌日期间记录到的“黑滴效应”} \label{fig_Venus_24}
\end{figure}
\subsubsection{二十世纪早期的进展}
在二十世纪之前，人们对金星几乎没有新的认识。其看似毫无特征的圆盘状外观未能透露其地表状况，直到光谱学和紫外观测的发展才揭示了更多秘密。

最早的紫外观测于 1920 年代进行，当时弗兰克·E·罗斯发现紫外照片中呈现出大量在可见光和红外波段中不存在的细节。他认为这是由于一层致密的黄色低层大气，上方覆盖着高层卷云所致\(^\text{[226]}\)。

人们注意到金星的圆盘缺乏可辨识的扁率，暗示其自转较慢，一些天文学家据此认为金星像当时被认为潮汐锁定的水星一样被潮汐锁定；但另一些研究者探测到金星夜侧存在大量热量，这又暗示其自转速度很快（当时人们尚未意识到其极高的表面温度），使问题更加复杂\(^\text{[227]}\)。后来在 1950 年代的研究表明，金星的自转为逆行。
\subsubsection{首次金星探测任务}
1961 年，人类首次尝试进行行星际空间飞行，即苏联“金星”计划中的无人探测器“金星 1 号”飞向金星，但在途中失去联系\(^\text{[228]}\)。

第一次成功的行星际任务，同样是前往金星，是美国“水手”计划的“水手 2 号”，于 1962 年 12 月 14 日以距金星表面 34,833\,km（21,644\,mi）的距离飞掠，并获取了金星大气的数据\(^\text{[229][230]}\)。

此外，最早的金星雷达观测在 1960 年代实施，并首次测得了其自转周期，其数值接近真实值\(^\text{[231]}\)。

1966 年发射的“金星 3 号”成为人类首个抵达并撞击月球以外其他天体的探测器和着陆器，但因坠毁于金星表面而未能返回数据。1967 年发射的“金星 4 号”成功在金星大气中实施科学实验后着陆。金星 4 号显示，金星表面温度几乎达 500\,°C（932\,°F），比水手 2 号计算的更高；其测得大气成分为 95\% 的二氧化碳（$CO_{2}$），并发现金星大气密度远超设计者的预期\(^\text{[232][233]}\)。

作为早期太空合作的范例，金星 4 号的数据与 1967 年的水手 5 号数据相结合，并由苏美联合科学团队在接下来一年的一系列专题会议中共同分析\(^\text{[234]}\)。

1970 年 12 月 15 日，“金星 7 号”成为首个在另一颗行星上软着陆并将数据传回地球的航天器\(^\text{[235]}\)。

1974 年，“水手 10 号”飞掠金星，以借助金星引力改变航向飞往水星，并拍摄了金星云层的紫外影像，揭示了金星大气中极高的风速。这是历史上首次使用的行星际引力助推技术，后来被众多探测器沿用。

1970 年代的雷达观测首次揭示了金星表面的地形细节。研究人员利用阿雷西博天文台 300\,m（1000\,ft）射电望远镜向金星发射无线电波，回波显示出两个反射率极高的区域，被命名为 Alpha 区和 Beta 区。观测还发现一个亮区，被归因于山地，并被命名为麦克斯韦山脉（Maxwell Montes）\(^\text{[236]}\)。这三处特征是目前金星上仅有的以非女性名称命名的地表特征\(^\text{[50]}\)。
\begin{figure}[ht]
\centering
\includegraphics[width=8cm]{./figures/9b4b2711eb58049a.png}
\caption{金星表面的首次影像及首次清晰的 180 度全景图，也是除地球外其他行星的首张此类全景图（1975 年，苏联“金星 9 号”着陆器）。黑白图像中呈现荒芜、黑色、板岩状的岩石，背景是一片平坦的天空。地面与探测器为画面主体。} \label{fig_Venus_26}
\end{figure}
1975 年，苏联的“金星 9 号”和“金星 10 号”着陆器传回了来自金星表面的首批图像，这些图像为黑白影像。美国宇航局随后通过“先驱者金星”计划获取了更多数据，该计划由两个独立任务组成\(^\text{[237]}\)：“先驱者金星多探测器”与“先驱者金星轨道器”，二者在 1978 年至 1992 年间绕金星运行\(^\text{[238]}\)。1982 年，苏联“金星 13 号”和“金星 14 号”着陆器首次获取了金星表面的单色彩色滤光片图像。在 1983 至 1984 年间，“金星 15 号”和“金星 16 号”在金星轨道运行，对其 25\% 的地表（从北极至北纬 30°）开展了详细测绘，之后苏联的“金星”探测计划宣告结束\(^\text{[239]}\)。
\begin{figure}[ht]
\centering
\includegraphics[width=6cm]{./figures/2cf993976ba79b80.png}
\caption{史密森学会乌德瓦－哈兹中心展出的“织女号”气球探测器} \label{fig_Venus_27}
\end{figure}
1985 年，苏联的“织女号”计划（Vega 1 与 Vega 2）携带了最后一批进入式探测器，并首次携带了地球以外的空气机器人（aerobots），通过充气气球的方式首次实现了地外大气中的飞行。

1990 至 1994 年间，“麦哲伦号”在金星轨道运行直至轨道衰减，并对金星表面进行了测绘。此外，“伽利略号”（1990）\(^\text{[240]}\)与“卡西尼–惠更斯号”（1998/1999）等探测器在飞往其他目的地途中也曾飞掠金星。
\subsubsection{重新开展的探索}
2006 年 4 月，欧洲航天局（ESA）的首次金星专门探测任务“金星快车”进入金星轨道。“金星快车”为金星大气带来了前所未有的观测成果。ESA 于 2014 年 12 月结束该任务，并于 2015 年 1 月使其再入\(^\text{[241]}\)。同年及次年，“信使号”（MESSENGER）在前往其他目的地途中也进行了金星飞掠。

2010 年，首艘成功的行星际太阳帆飞船 IKAROS 前往金星并实施了飞掠。

2015 至 2024 年间，日本的“曙光号”（Akatsuki）在金星轨道上开展探测，同时“贝皮科伦布号”（BepiColombo）在 2020/2021 年进行了飞掠。
\begin{figure}[ht]
\centering
\includegraphics[width=8cm]{./figures/51a8621358abd9f4.png}
\caption{于 2021 年拍摄了这段金星夜侧的可见光影像，透过云层显示出其炽热且微弱发光的表面，以及作为大片暗区的阿芙洛狄忒高地；而在白昼侧，由于云层被照亮，此类观测无法进行\(^\text{[242][243]}\)。} \label{fig_Venus_28}
\end{figure}
\subsubsection{正在执行与计划中的任务}
\begin{figure}[ht]
\centering
\includegraphics[width=8cm]{./figures/25fe75bbacd989dd.png}
\caption{金星全球地形图，并标注所有探测器着陆点（红色：返回图像；黑点：采集样本并进行现场分析）} \label{fig_Venus_29}
\end{figure}
截至 2025 年，金星附近没有正在运行的探测器，帕克太阳探测器计划在 2030 年之前多次返回金星进行飞掠。

多项探测器正在研制中，同时还有多项任务仍处于早期概念阶段。下一项已计划的金星任务是“金星生命探测器”（Venus Life Finder），预计最早于 2026 年夏季发射。

印度空间研究组织（ISRO）正在推进金星轨道器任务（Venus Orbiter Mission），计划于 2028 年发射。阿联酋的 MBR 探索者（MBR Explorer）小行星任务将对金星进行一次飞掠。美国宇航局（NASA）已批准两项金星任务——VERITAS 与 DAVINCI——计划最早于 2031 年发射。欧洲航天局（ESA）计划于 2031 年发射 EnVision。

\textbf{目标}

金星被确定为未来研究中的重要对象，用于理解：
\begin{itemize}
\item 太阳系与地球的起源，以及类似我们这样的行星系统在宇宙中是普遍还是稀有；
\item 行星天体如何从原始状态演化至今日多样的形态；
\item 导致适宜生存环境与生命出现的条件如何发展[244]。
\end{itemize}
\subsubsection{载人任务概念}
自 1960 年代起，金星便被视作火星载人探测的潜在中转点，即采用“冲日任务”，而非直接利用金星引力辅助的“合日任务”。研究显示，这类方案使前往火星的任务更快、更安全，并拥有更好的返程或中止窗口，且辐射暴露不高于或甚至低于直接前往火星的任务\(^\text{[245][246]}\)。

\textbf{金星大气层中的可能居住}
\begin{figure}[ht]
\centering
\includegraphics[width=8cm]{./figures/2019ebdd05da8040.png}
\caption{美国宇航局“金星高空作业构想”（HAVOC）载人浮空前哨基地的艺术示意图} \label{fig_Venus_30}
\end{figure}
尽管金星表面的环境极端恶劣，但在距地表 50\,km 的高度，其大气压、温度以及所受的太阳与宇宙辐射均与地球表面相似（“温和条件”）\(^\text{[247][114][113][183]}\)。然而，若要在金星大气中建立任何形式的人类存在，面临的众多工程挑战之一是其大气中具有强腐蚀性的高浓度硫酸\(^\text{[248]}\)。作为替代在火星等行星表面生活的流行设想，人们提出了用于载人探索甚至永久“浮空城市”的金星大气层浮空器方案\(^\text{[248][249][250][251][252]}\)。美国宇航局的“金星高空作业构想”（HAVOC）则是一项用于研究载人浮空器设计的训练性概念。
\subsection{生命的可能性}
自 20 世纪 60 年代初以来，关于金星表面存在生命的推测显著减少，因为人们逐渐认识到金星表面的环境相比地球极端得多。金星的极端温度与大气压使目前已知的以水为基础的生命形式难以存在。

一些科学家推测，在金星较高、较冷且酸性的上层大气中，可能存在嗜热嗜酸的极端微生物\(^\text{[253][254][255]}\)。这类推测可追溯到 1967 年，当时卡尔·萨根和哈罗德·J·莫罗维茨在《自然》杂志发文指出，金星云层中探测到的微小颗粒或许是类似地球细菌的生物（其尺度相似）：

“虽然金星地表条件使生命假说不太可信，但金星云层则是完全不同的故事。正如数年前有人指出的那样，水、二氧化碳和阳光——光合作用所需的要素——在云层附近极为丰富。”\(^\text{[256]}\)

2019 年 8 月，由李延珠（Yeon Joo Lee）领导的天文学团队报告称，金星大气中“未知吸收体”引起了长期的吸收率与反照率变化，这些吸收体可能是化学物质，也可能是在高空存在的大型微生物群落，它们会影响金星的气候\(^\text{[137]}\)。其光吸收曲线与地球云层中的微生物几乎完全一致。其他研究也得出了类似结论\(^\text{[257]}\)。

2020 年 9 月，由卡迪夫大学的简·格里夫斯（Jane Greaves）领导的天文学团队宣布，他们可能在金星上层云层中探测到磷化氢，这是一种已知不会由金星地表或大气中的化学过程产生的气体\(^\text{[258][116][115][259][260]}\)。一种可能的解释是生命活动\(^\text{[261]}\)。该磷化氢在距地表至少 48\,km（30\,mi）高度、主要位于中纬度地区被探测到，而在两极则没有发现。这一发现促使 NASA 管理员吉姆·布里登斯廷公开呼吁重新关注金星研究，称磷化氢的发现是“在人类寻找地外生命道路上最重要的发展之一”\(^\text{[262][263]}\)。

随后对用于识别金星大气中磷化氢的数据处理方式的分析提出了质疑：用于拟合的 12 阶多项式可能放大了噪声并生成了假信号（参见龙格现象）。在其他电磁波段中对金星大气的观测未能检出预期的磷化氢吸收线\(^\text{[264]}\)。到 2020 年 10 月下旬，重新分析并正确扣除背景后，已不再显示具有统计显著性的磷化氢信号\(^\text{[265][266][267]}\)。

格里夫斯团队的部分成员正在参与麻省理工学院的计划，准备与 Rocket Lab 火箭公司合作发射首个私人行星际探测器“金星生命探测器”（Venus Life Finder），通过探测器进入金星大气以寻找有机物\(^\text{[269]}\)。
\subsubsection{行星保护}
由于金星表面环境极其严酷，金星被归类为行星保护第二类（category two），即第二低级别\(^\text{[269]}\)。这意味着航天器携带的星际污染几乎不可能影响对其的科学研究。

然而，随着金星可能存在生物标志物的发现，对至少部分金星大气层的分类提出了质疑。但由于这些层尚未被确认为能够支持生命，因此不建议调整其行星保护分类\(^\text{[270]}\)。
\subsubsection{人类可居性}
尽管金星表面对人类极不友好，但在距地表 50\,km 的高度，其大气条件不仅被认为可能适宜本地生命，也可能适宜人类生存，比太阳系除地球以外的任何地方都更具优势。从大气压、重力、温度到辐射，除了化学成分外，这里的条件几乎与地球表面相同。基于这一前景，有人提出可在此高度建造人类的浮空居住设施，以实现人类前往金星的可能性\(^\text{[271]}\)。
\subsection{文化中的金星}
金星是夜空中最显著的天体之一，在众多文化的神话、占星学与文学作品中都被赋予了特殊的重要地位。
\begin{figure}[ht]
\centering
\includegraphics[width=8cm]{./figures/a5aa0ac67bd47cec.png}
\caption{} \label{fig_Venus_31}
\end{figure}
多首赞美诗歌颂伊南娜作为金星女神的角色\(^\text{[205][272][273]}\)。神学教授杰弗里·库利（Jeffrey Cooley）认为，在许多神话中，伊南娜的行动可能与金星在天空中的运行相对应\(^\text{[205]}\)。金星不连续的运动与伊南娜的神话形象及其双重性相关\(^\text{[205]}\)。在《伊南娜降入冥界》中，伊南娜能够降入冥界又重返天界，这是其他神祇无法做到的。金星似乎也有类似的“降落”，它在西方落下，又在东方升起\(^\text{[205]}\)。一首引言性赞美诗描述伊南娜离开天界前往库尔（Kur），被认为是山地，其意象与伊南娜在西方的升落相对应\(^\text{[205]}\)。《伊南娜与舒卡勒图达》《伊南娜降入冥界》中都出现了与金星运行相平行的描绘\(^\text{[205]}\)。《伊南娜与舒卡勒图达》中，舒卡勒图达被描述为在天穹中寻找伊南娜，可能指在东方与西方地平线上寻找\(^\text{[274]}\)。在同一神话中，伊南娜在寻找攻击者时的多次移动也与金星在天空中的运行一致\(^\text{[205]}\)。

通过美索不达米亚文化的影响，古埃及人与希腊人可能在公元前二千纪（或最晚在更后的晚期）便知道晨星与昏星为同一天体\(^\text{[275][276]}\)。埃及人称晨星为 Tioumoutiri，昏星为 Ouaiti\(^\text{[277]}\)。他们起初将金星描绘为凤凰或苍鹭（见 Bennu）\(^\text{[275]}\)，称其为“穿越者”或“带十字之星”\(^\text{[275]}\)，并与奥西里斯相联系；后又将其描绘为双头形象（人首或隼首），并与荷鲁斯相关联\(^\text{[276]}\)，荷鲁斯是伊西斯之子（在更后的希腊化时期，伊西斯与哈索尔被认同为阿佛洛狄忒）。希腊人将晨星称为 Phōsphoros（光明使者，磷元素名“phosphorus”即源于此；或 Ēōsphoros，意为“黎明带来者”），将昏星称为 Hesperos（意为“西方者”）\(^\text{[278]}\)，二者皆为黎明女神 Eos 之子，因此是阿佛洛狄忒的外孙。尽管到罗马时代二者已被认作同一天体并统称为“维纳斯之星”，但传统的两个希腊名称仍然沿用，并通常以拉丁语译为 Lucifer 与 Vesper\(^\text{[278][279]}\)。

古典诗人如荷马、萨福、奥维德与维吉尔都描写过金星及其光辉\(^\text{[280]}\)。诗人威廉·布莱克、罗伯特·弗罗斯特、莱蒂希娅·伊丽莎白·兰登、阿尔弗雷德·丁尼生与威廉·华兹华斯也都为金星写过颂歌\(^\text{[281]}\)。作曲家霍尔斯特在其《行星组曲》中将金星作为第二乐章。

在印度，金星的梵文名称为 Shukra Graha，意为“圣者 Shukra 之星”，Shukra 是一位强大的圣人。在吠陀占星学中\(^\text{[282]}\)，Shukra 的梵文含义为“清澈、纯净”或“明亮、清晰”。作为九曜之一，它被认为影响财富、享乐与繁衍；它是 Bhrgu 之子，Daityas 的导师，也是 Asuras 的古鲁\(^\text{[283]}\)。

英语名称“Venus”源自古罗马人。罗马人以爱神维纳斯为该行星命名，而维纳斯的原型是希腊的爱神阿佛洛狄忒\(^\text{[284]}\)，而阿佛洛狄忒又源自苏美尔宗教中的伊南娜（阿卡德宗教中称为伊什塔尔），她们皆与金星相关\(^\text{[273][272]}\)。金星与这些女神对应的星期为星期五（Friday），以日耳曼女神 Frigg 命名，而 Frigg 也与罗马的维纳斯相关。

中文中，金星称为“金星”，意为“金之星”；在中国传统五行中，金星属金。这一传统在现代中、日、韩、越文化中共享，各语言中金星均字面意为“金星”（金星）\(^\text{[285][286][287][288]}\)。

玛雅人认为金星仅次于太阳和月亮，是最重要的天体。他们称金星为 Chac ek\(^\text{[289]}\)或 Noh Ek'（意为“大星”）\(^\text{[290]}\)。金星的周期对其历法十分重要，并在他们的一些典籍中记载，例如《墨西哥玛雅手抄本》和《德累斯顿抄本》。智利国旗（孤星旗）所描绘的“孤星”即象征金星。
\subsubsection{现代文化}
\begin{figure}[ht]
\centering
\includegraphics[width=8cm]{./figures/47a31ae5a698ebd5.png}
\caption{在文森特·梵高 1889 年的绘画《星夜》中，金星被描绘在画面中那棵巨大柏树的右侧\(^\text{[291][292]}\)。} \label{fig_Venus_32}
\end{figure}
无法穿透的金星云层使科幻作家得以自由想象其表面状况；尤其在早期观测显示金星不仅与地球大小相近，还拥有稠密大气之后，他们更是大胆发挥。由于金星比地球更靠近太阳，它常被描绘为更温暖，但仍适宜人类居住\(^\text{[293]}\)。这一文学类型在 1930 至 1950 年代达到巅峰，当时科学已揭示金星的一些特征，但尚不了解其表面的严酷现实。首次金星探测任务的发现证明真实情况截然不同，从而终结了这一体裁\(^\text{[294]}\)。随着对金星科学认知的进步，科幻作家也试图同步发展，尤其通过描绘人类试图地球化（金星改造）金星的设想\(^\text{[295]}\)。
\subsubsection{符号}
圆圈下带一小十字的符号即所谓的金星符号，因作为金星的天文符号而得名。该符号源于古希腊，更广泛地代表女性特质，并被生物学采用为女性性别符号\(^\text{[296][297][298]}\)，与用于男性的火星符号相对应，有时也与用作雌雄同体现象的水星符号并列。这种将金星与火星赋以性别含义的做法，也被用于异性恋规范化的配对语境中，用以将女性与男性刻板地描述为来自完全不同的“行星”。这种理解在 1992 年出版的《男人来自火星，女人来自金星》一书中被进一步普及\(^\text{[299]}\)。

金星符号也曾用于西方炼金术中，代表元素铜（如同水星符号亦代表元素汞）\(^\text{[297][298]}\)。由于抛光铜自古以来一直用于制作镜子，金星符号有时被称为“金星之镜”，象征爱神的镜子，尽管这一解释后来被认为不太可能是其真正来源\(^\text{[297][298]}\)。

除金星符号外，还有许多其他符号与金星相关，其中常见的还有新月，或尤以星形符号著称，如伊什塔尔之星\(^\text{[300]}\)。
\subsection{另见}
\begin{itemize}
\item 金星纲要  
\item 太阳系行星的物理性质  
\item 金星带 
\end{itemize} 
\subsection{注释 } 
\begin{enumerate}
\item 在新闻稿与科学出版物中被误写为 “Ganiki Chasma”\(^\text{[66]}\)。  
\item 必须明确“接近度”一词的含义。在天文学文献中，“最近行星”（closest planets）常指两颗行星在其轨道上彼此最接近的程度，即两者轨道之间的最小距离。然而，这并不意味着这两颗行星在时间尺度上总体上彼此最近。由于水星比金星更靠近太阳，水星在更长时间内与地球保持相对接近；因此，从时间平均的角度说，可以认为“水星是与地球平均意义上最近的行星”。然而，在采用这种时间平均定义时，水星实际上是整个太阳系中与所有其他行星平均距离最近的行星。因此，这种基于平均距离的“接近”定义似乎并不特别有用。BBC Radio 4 节目《More or Less》的某一集对这些不同的“接近”概念进行了相当清晰的解释\(^\text{[173]}\)。  
\item 若干中世纪伊斯兰天文学家声称的凌日观测后来被证明是太阳黑子[217]。阿维森纳并未记录他的观测日期。在他的一生中确实发生过一次金星凌日，即 1032 年 5 月 24 日，但尚不确定当时从他的所在地是否能够看到这一现象\(^\text{[218]}\)。
\end{enumerate}
\subsection{参考文献}
\begin{enumerate}
\item Ricardo, Nunes. “Ricardo Nunes Astronomy Page”。[www.astrosurf.com。检索于](http://www.astrosurf.com。检索于) 2025 年 11 月 19 日。
\item “Venusian”。Lexico 英国英语词典。牛津大学出版社。2020 年 3 月 23 日存档。
\item “Venusian”。Merriam-Webster.com 词典。Merriam-Webster。
\item “Cytherean”。《牛津英语词典》（在线版）。牛津大学出版社。（需订阅或机构访问权限。）
\item “Venerean, Venerian”。《牛津英语词典》（在线版）。牛津大学出版社。（需订阅或机构访问权限。）
\item Williams, David R.（2020 年 11 月 25 日）。“Venus Fact Sheet”。美国宇航局戈达德太空飞行中心。2018 年 5 月 11 日存档。2021 年 4 月 15 日检索。
\item Yeomans, Donald K. “Horizons Web-Interface for Venus (Major Body=2)”。JPL Horizons 在线星历系统。2023 年 8 月 18 日存档。2010 年 11 月 30 日检索。——选择“Ephemeris Type: Orbital Elements”“Time Span: 2000-01-01 12:00 to 2000-01-02”。（“Target Body: Venus”“Center: Sun”保持默认。）结果为 J2000 历元的瞬时摄动轨道元素。
\item Simon, J.L.; Bretagnon, P.; Chapront, J.; Chapront-Touzé, M.; Francou, G.; Laskar, J.（1994 年 2 月）。“月球与行星的岁差公式与平均轨道要素的数值表达式”。《天文学与天体物理学》282（2）：663–683。Bibcode:1994A&A...282..663S。
\item Souami, D.; Souchay, J.（2012 年 7 月）。“太阳系的不变量面”。《天文学与天体物理学》543：11。Bibcode:2012A&A...543A.133S。doi:10.1051/0004-6361/201219011。A133。
\item Seidelmann, P. Kenneth; Archinal, Brent A.; A'Hearn, Michael F.; 等（2007）。“IAU/IAG 工作组关于行星和卫星制图坐标和自转参量（2006）的报告”。《天体力学与动力天文学》98（3）：155–180。Bibcode:2007CeMDA..98..155S。doi:10.1007/s10569-007-9072-y。
\item Konopliv, A. S.; Banerdt, W. B.; Sjogren, W. L.（1999 年 5 月）。“金星引力：180 阶度模型”（PDF）。《伊卡洛斯》139（1）：3–18。Bibcode:1999Icar..139....3K。CiteSeerX 10.1.1.524.5176。doi:10.1006/icar.1999.6086。2010 年 5 月 26 日从原 PDF 存档。
\item “行星与冥王星：物理特征”。NASA。2008 年 11 月 5 日。2006 年 9 月 7 日存档。2015 年 8 月 26 日检索。
\item “行星事实”。行星学会（The Planetary Society）。2012 年 5 月 11 日存档。2016 年 1 月 20 日检索。
\item Margot, Jean-Luc; Campbell, Donald B.; Giorgini, Jon D.; 等（2021 年 4 月 29 日）。“金星的自旋状态与转动惯量”。《自然·天文学》（Nature Astronomy）。5（7）：676–683。arXiv:2103.01504。Bibcode:2021NatAs...5..676M。doi:10.1038/s41550-021-01339-7。S2CID 232092194。
\item “IAU/IAG 行星与卫星制图坐标和自转要素工作组报告”。国际天文学联合会。2000 年。2020 年 5 月 12 日存档。2007 年 4 月 12 日检索。
\item Archinal, B. A.; Acton, C. H.; A'Hearn, M. F.; Conrad, A.; Consolmagno, G. J.; Duxbury, T.; Hestroffer, D.; Hilton, J. L.; Kirk, R. L.; Klioner, S. A.; McCarthy, D.; Meech, K.; Oberst, J.; Ping, J.; Seidelmann, P. K.（2018）。“IAU 制图坐标与自转要素工作组报告：2015”。《天体力学与动力天文学》。130（3）：22。Bibcode:2018CeMDA.130...22A。doi:10.1007/s10569-017-9805-5。ISSN 0923-2958。
\item Mallama, Anthony; Krobusek, Bruce; Pavlov, Hristo（2017）。“行星的宽波段星等与反照率的综合研究，并应用于系外行星和第九行星假说”。《伊卡洛斯》（Icarus）。282：19–33。arXiv:1609.05048。Bibcode:2017Icar..282...19M。doi:10.1016/j.icarus.2016.09.023。S2CID 119307693。
\item Haus, R.; Kappel, D.; Arnoldb, G.（2016 年 7 月）。“基于改进的金星中低层大气模型的辐射能量平衡”（PDF）。《伊卡洛斯》（Icarus）。272：178–205。Bibcode:2016Icar..272..178H。doi:10.1016/j.icarus.2016.02.048。2017 年 9 月 22 日 PDF 存档。2019 年 6 月 25 日检索。
\item “大气与行星温度”。美国化学学会。2013 年 7 月 18 日。2023 年 1 月 27 日存档。2023 年 1 月 3 日检索。Herbst, K.; Banjac, S.; Atri D.; Nordheim, T. A（2020 年 1 月 1 日）。“在宜居性背景下重新评估宇宙射线导致的金星辐射剂量”。《天文学与天体物理学》（Astronomy & Astrophysics）。633。图 6。arXiv:1911.12788。Bibcode:2020A&A...633A..15H。doi:10.1051/0004-6361/201936968。ISSN 0004-6361。S2CID 208513344。2021 年 12 月 5 日存档。2021 年 11 月 20 日检索。
\item Mallama, Anthony; Hilton, James L.（2018 年 10 月）。“为《天文年鉴》计算行星视星等”。《天文学与计算》（Astronomy and Computing）。25：10–24。arXiv:1808.01973。Bibcode:2018A&C....25...10M。doi:10.1016/j.ascom.2018.08.002。S2CID 69912809。
\item “百科全书——最亮的天体”。IMCCE。2023 年 7 月 24 日存档。2023 年 5 月 29 日检索。
\item Lawrence, Pete（2005）。《寻找金星之影》。Digitalsky.org.uk。2012 年 6 月 11 日存档。2012 年 6 月 13 日检索。
\item Walker, John。《在大白天观测金星》。Fourmilab Switzerland。2017 年 3 月 29 日存档。2017 年 4 月 19 日检索。
\item Andrews, Roy George（2024 年 5 月 27 日）。《金星上的熔岩河揭示更为活跃的火山行星——新软件让科学家重新分析旧雷达影像，提供了火山仍在重塑这颗地狱般行星的最强证据之一》。《纽约时报》。2024 年 5 月 28 日存档。2024 年 5 月 28 日检索。
\item Sulcanese, Davide; Mitri, Giuseppe; Mastrogiuseppe, Marco（2024 年 5 月 27 日）。《麦哲伦雷达揭示金星正在进行的火山活动证据》。《自然·天文学》（Nature Astronomy）。8（8）：973–982。Bibcode:2024NatAs...8..973S。doi:10.1038/s41550-024-02272-1。ISSN 2397-3366。2024 年 5 月 28 日存档。2024 年 5 月 28 日检索。
\item Chang, Kenneth（2023 年 10 月 26 日）。《数十亿年前，金星或曾拥有关键的类地特征——新研究提出，这颗地狱般的第二行星可能曾有板块构造，使其更适宜生命》。《纽约时报》。2023 年 10 月 26 日存档。2023 年 10 月 27 日检索。
\item Weller, Matthew B.; 等（2023 年 10 月 26 日）。《金星大气中的氮由古代板块构造解释》。《自然·天文学》（Nature Astronomy）。7（12）：1436–1444。Bibcode:2023NatAs...7.1436W。doi:10.1038/s41550-023-02102-w。S2CID 264530764。2023 年 10 月 27 日存档。2023 年 10 月 27 日检索。
\item Jakosky, Bruce M.（1999）。《类地行星的大气》。载于 Beatty, J. Kelly；Petersen, Carolyn Collins；Chaikin, Andrew（编）。《新太阳系》（第 4 版）。波士顿：Sky Publishing。第 175–200 页。ISBN 978-0-933346-86-4。OCLC 39464951。
\item Hashimoto, George L.; Roos-Serote, Maarten; Sugita, Seiji; Gilmore, Martha S.; Kamp, Lucas W.; Carlson, Robert W.; Baines, Kevin H.（2008 年 12 月 31 日）。《伽利略近红外成像光谱仪数据暗示金星具有长英质高地地壳》。《地球物理研究杂志：行星》（Journal of Geophysical Research: Planets）。113（E5），2008JE003134。《地球与空间科学进展》。Bibcode:2008JGRE..113.0B24H。doi:10.1029/2008JE003134。S2CID 45474562。
\item Shiga, David（2007 年 10 月 10 日）。《金星的古老海洋是否孕育了生命？》。《新科学家》。2009 年 3 月 24 日存档。2017 年 9 月 17 日检索。
\item Lopes, Rosaly M. C.; Gregg, Tracy K. P.（2004）。《火山世界：探索太阳系火山》。施普林格出版社。第 61 页。ISBN 978-3-540-00431-8。
\item Squyres, Steven W.（2016）。“金星”。《大英百科全书在线》。2014 年 4 月 28 日存档。2016 年 1 月 7 日检索。
\item Darling, David。《金星》。《科学百科全书》。英国邓迪。2021 年 10 月 31 日存档。2022 年 3 月 24 日检索。
\item Lammer, Helmut; Brasser, Ramon; Johansen, Anders; Scherf, Manuel; Leitzinger, Martin（2021 年 2 月）。《金星、地球与火星的形成：由同位素约束》。《空间科学评论》（Space Science Reviews）。217（1）：7。arXiv:2102.06173。doi:10.1007/s11214-020-00778-4。ISSN 0038-6308。
\item “第 6 章：迈向巨型化——NASA 科学”。2018 年 12 月 11 日。2025 年 9 月 7 日检索。
\item Koren, Marina（2022 年 2 月 3 日）。《我们的太阳系以真实色彩呈现时竟如此不同凡响》。《大西洋月刊》。2025 年 6 月 24 日检索。
\item Mueller, Nils（2014）。《金星的表面与内部》。见 Tilman Spohn；Doris Breuer；T. V. Johnson（编），《太阳系百科全书》（第 3 版）。牛津：爱思唯尔科学与技术出版社。ISBN 978-0-12-415845-0。2021 年 9 月 29 日存档。2016 年 1 月 12 日检索。
\item Esposito, Larry W.（1984 年 3 月 9 日）。《二氧化硫：间歇性喷注显示金星存在活跃火山作用的证据》。《科学》（Science）。223（4640）：1072–1074。Bibcode:1984Sci...223.1072E。doi:10.1126/science.223.4640.1072。PMID 17830154。S2CID 12832924。
\item Bullock, Mark A.; Grinspoon, David H.（2001 年 3 月）。《金星气候的近期演化》（PDF）。《伊卡洛斯》（Icarus）。150（1）：19–37。Bibcode:2001Icar..150...19B。CiteSeerX 10.1.1.22.6440。doi:10.1006/icar.2000.6570。2003 年 10 月 23 日存档（PDF）。
\item Basilevsky, Alexander T.; Head, James W. III（1995）。《金星的全球地层学：对三十六个随机测试区域的分析》。《地球、月球与行星》（Earth, Moon, and Planets）。66（3）：285–336。Bibcode:1995EM&P...66..285B。doi:10.1007/BF00579467。S2CID 21736261。
\item Jones, Tom；Stofan, Ellen（2008）。《行星学：解锁太阳系的奥秘》。国家地理学会。第 74 页。ISBN 978-1-4262-0121-9。2017 年 7 月 16 日存档。2017 年 4 月 20 日检索。
\item Kaufmann, W. J.（1994）。《宇宙》。纽约：W. H. Freeman。第 204 页。ISBN 978-0-7167-2379-0。
\item 《国家地理》（2024）。《金星具有活跃火山生命》。
\item Nimmo, F.; McKenzie, D.（1998）。《金星的火山作用与构造》。《年度地球与行星科学评论》。26（1）：23–53。Bibcode:1998AREPS..26...23N。doi:10.1146/annurev.earth.26.1.23。S2CID 862354。
\item Strom, Robert G.; Schaber, Gerald G.; Dawson, Douglas D.（1994 年 5 月 25 日）。《金星的全球再表面化》。《地球物理研究杂志》。99（E5）：10899–10926。Bibcode:1994JGR....9910899S。doi:10.1029/94JE00388。2020 年 9 月 16 日存档。2019 年 6 月 25 日检索。
\item Frankel, Charles（1996）。《太阳系的火山》。剑桥大学出版社。ISBN 978-0-521-47770-3。2023 年 1 月 30 日检索。
\item Mitchell, Don P.（1975 年 10 月 20 日）。《苏联金星图像》。Don P. Mitchell 网站。2025 年 6 月 2 日检索。
\item Moroz, V. I.; Golovin, Yu. M.; Ekonomov, A. P.; Moshkin, B. E.; Parfent'ev, N. A.; San'ko, N. F.（1980）。《金星白天天空的光谱》。《自然》（Nature）。284（5753）：243–244。Bibcode:1980Natur.284..243M。doi:10.1038/284243a0。ISSN 0028-0836。2025 年 6 月 2 日检索。
\item Batson, R. M.; Russell, J. F.（1991 年 3 月 18–22 日）。《为金星新发现的地貌命名》（PDF）。《第 22 届月球与行星科学会议论文集》。美国德克萨斯州休斯敦。第 65 页。Bibcode:1991pggp.rept..490B。2011 年 5 月 13 日存档（PDF）。2009 年 7 月 12 日检索。
\item Young, Carolynn（编）（1990 年 8 月 1 日）。《麦哲伦金星探测指南》。加利福尼亚：喷气推进实验室。第 93 页。2016 年 12 月 4 日存档。2016 年 1 月 13 日检索。
\item Davies, M. E.; Abalakin, V. K.; Bursa, M.; Lieske, J. H.; Morando, B.; Morrison, D.; Seidelmann, P. K.; Sinclair, A. T.; Yallop, B.; Tjuflin, Y. S.（1994）。《国际天文学联合会工作组关于行星和卫星的制图坐标及自转要素的报告》。《天体力学与动力天文学》。63（2）：127–148。Bibcode:1996CeMDA..63..127D。doi:10.1007/BF00693410。S2CID 189850694。
\item Young, Carolynn（编）（1990 年 8 月 1 日）。《麦哲伦金星探测指南》。加利福尼亚：喷气推进实验室。第 99–100 页。2016 年 12 月 4 日存档。2016 年 1 月 13 日检索。
\item Helbert, Jörn；Müller, Nils；Kostama, Petri；Marinangeli, Lucia；Piccioni, Giuseppe；Drossart, Pierre（2008）。《金星快车号 VIRTIS 观测到的拉达地地区表面亮度变化及其对区域演化的启示》（PDF）。《地球物理研究通讯》。35（11）：L11201。Bibcode:2008GeoRL..3511201H。doi:10.1029/2008GL033609。ISSN 1944-8007。
\item Petkowski, Janusz；Seager, Sara（2021 年 11 月 18 日）。《金星曾经拥有过海洋吗？》*。MIT“金星云层生命”项目。2023 年 4 月 13 日存档。2023 年 4 月 13 日检索。
\item Gilmore, Martha；Treiman, Allan；Helbert, Jörn；Smrekar, Suzanne（2017 年 11 月 1 日）。《通过观测与实验约束的金星表面成分》。《空间科学评论》。212（3）：1511–1540。Bibcode:2017SSRv..212.1511G。doi:10.1007/s11214-017-0370-8。ISSN 1572-9672。S2CID 126225959。
\item 《新的目录在金星最优质的地表图像中精确定位了火山锥——这些图像由 NASA 的麦哲伦号在 30 年前获取》。skyandtelescope.org。2023 年 4 月 14 日。2023 年 4 月 16 日存档。2023 年 4 月 16 日检索。
\item Hahn, Rebecca M.; Byrne, Paul K.（2023 年 4 月）。《金星火山的形态学与空间分布分析》。《地球物理研究期刊：行星》。128（4）：e2023JE007753。Bibcode:2023JGRE..12807753H。doi:10.1029/2023JE007753。S2CID 257745255。作者利用麦哲伦号合成孔径雷达左右视全球拼接图（75 米/像素），建立了一份包含约 85,000 座火山的全球目录，其中约 99\% 的火山直径小于 5 公里。研究发现金星拥有远多于以往绘制的火山，虽然它们几乎遍布整个星球，但大小–频率分布分析显示直径 20–100 公里范围内的火山数量相对不足，这可能与岩浆供应和喷发速率有关。
\item Karttunen, Hannu；Kroger, P.；Oja, H.；Poutanen, M.；Donner, K. J.（2007）。《基础天文学》。施普林格出版社。第 162 页。ISBN 978-3-540-34143-7。2023 年 1 月 30 日检索。
\item Bauer, Markus（2012 年 12 月 3 日）。《金星火山是否被当场“抓包”？》。欧洲航天局。2021 年 4 月 14 日存档。2021 年 4 月 14 日检索。
\item Glaze, Lori S.（1999 年 8 月）。《金星上爆炸性火山活动对 SO₂ 的输运作用》。《地球物理研究杂志》。104（E8）：18899–18906。Bibcode:1999JGR...10418899G。doi:10.1029/1998JE000619。
\item Marcq，Emmanuel; Bertaux，jean-loup; Montmessin，Franck; Belyaev，Denis (2013年1月)。“金星动态大气云顶二氧化硫的变化”。自然地球科学。6(1):25-28.Bibcode:2013NatGe...6...25M。doi:10.1038/ngeo1650。S2CID59323909。 
\item 霍尔，桑农 (2020年1月9日)。金星上的火山可能还在冒烟 -- 地球上的行星科学实验表明，太阳的第二颗行星可能正在进行火山活动。。纽约时报。已存档从原件到2020年1月9日。已检索1月10日2020。
\item 贾斯汀·菲利贝托 (2020年1月3日)。“从橄榄石的风化速率证明了当今金星上的火山活动”。科学。6(1) eaax7445。Bibcode:2020年SciA。6.7445F。doi:10.1126/sciadv.aax7445。PMC6941908。PMID31922004。  
\item “早期，充满活力的碰撞可能助长了金星火山活动: 研究 | Sci.News”。2023年7月20日已存档从2023年7月21日开始。已检索7月21日2023。
\item "Ganis Chasma"。行星命名法地名词典。USGS天体地质科学中心。已存档从原件到2018年10月13日。已检索4月14日2021。
\item 艾米丽·拉克达瓦拉(2015年6月18日)。“金星上的瞬态热点: 活火山活动的最佳证据”。行星社会。已存档从原件到2015年6月20日。已检索6月20日2015。
\item “在金星上发现的热熔岩流”。欧洲航天局。2015年6月18日存档自原来的2015年6月19日。已检索6月20日2015。
\item Shalygin, E.五、；Markiewicz, W.J.;Basilevsky, A.T.;蒂托夫，D.五、；Ignatiev, N.I.;头，J.W。(2015年6月17日)。“Ganiki Chasma裂谷带金星上的活火山活动”。地球物理研究快报。42(12):4762-4769。Bibcode:2015乔治。。42.4762S。doi:10.1002/2015GL064088。S2CID16309185。 
\item 克鲁格，杰弗里 (2023年3月17日)。“为什么在金星上发现活火山很重要”。时间。已存档从2023年3月19日开始。已检索3月19日2023。
\item 罗密欧，我；特科特，D。L.(2009)。“金星上火山单位的频率区域分布: 对行星表面重铺的影响”(PDF)。伊卡洛斯。203(1):13-19。Bibcode:2009年伊卡尔。。203。13R。doi:10.1016/j.icarus.2009.03.036。已存档(PDF)从2019年12月19日的原件开始。已检索12月15日2018。
\item 赫里克，R。R、；菲利普斯，R。J.(1993)。“金星大气对传入流星体和撞击坑人口的影响”。伊卡洛斯。112(1):253-281。Bibcode:1994年伊卡尔。。112。。253小时。doi:10.1006/icar.1994.1180。
\item 莫里森，大卫; 欧文斯，托比亚斯C。(2003)。行星系统(第三版)。旧金山:本杰明·卡明斯。ISBN 978-0-8053-8734-6。
\item Goettel, K.A.;盾牌，J.A.;德克尔，D。A.(1981年3月16日至20日)。“金星组成的密度限制”。月球与行星科学会议论文集。德克萨斯州休斯顿:佩加蒙出版社。pp。 1507-1516。Bibcode:1982LPSC...12.1507G。
\item 福雷，冈特; 梅辛，特蕾莎·M。(2007)。行星科学导论: 地质学视角。斯普林格电子书收藏。斯普林格。第201页。ISBN 978-1-4020-5233-0。
\item 杜穆林，C.；托比，G.；Verhoeven, O.;罗森布拉特，P.；Rambaux, N.(2017年6月)。“金星内部的潮汐限制”(PDF)。地球物理研究杂志: 行星。122(6):1338-1352。Bibcode:2017JGRE..122.1338D。doi:10.1002/2016JE005249。S2CID134766723。已存档(PDF)从原件到2020年5月9日。已检索5月3日2021。  
\item 艾塔，A。(2012年4月)。“金星的内部结构，温度和核心组成”。伊卡洛斯。218(2):967-974。Bibcode:2012Icar .. 218。。967A。doi:10.1016/j.icarus.2012.01.007。已存档从原件到2021年9月29日。已检索1月17日2016。
\item 乔纳森·奥卡拉汉 (2021年4月29日)。“我们首次测量了金星行星核心的大小。。新科学家。已存档从原始的2021年5月2日。已检索5月2日2021。
\item Semprich，Julia; Filiberto，Justin; Weller，Matthew; Gorce，Jennifer; Clark，Nolan (2025年3月25日)。“金星的变质作用是地壳厚度和再循环的驱动力”。自然通讯。16(1): 2905。Bibcode:2025年纳科。。16.2905s。doi:10.1038/s41467-025-58324-1。ISSN2041-1723。PMC11937330。PMID40133342。   
\item 尼莫，F。(2002)。“从麦哲伦卫星观测到的金星地壳分析阿塔兰塔Planitia，Beta Regio和thedis Regio”。地质学。30(11):987-990。Bibcode:2002Geo .... 30 .. 987N。doi:10.1130/0091-7613(2002)030<0987:WDVLAM>2.0.CO;2。ISSN0091-7613。S2CID13293506。  
\item 多尔吉诺夫，Sh.；埃罗申科，E.G.;刘易斯，L。(1969年9月)。“金星附近磁场的性质”。宇宙研究。7: 675。Bibcode:1969年科斯雷。7 .. 675D。
\item Kivelson, G.M、；罗素，C。T.(1995)。空间物理学导论。剑桥大学出版社。ISBN 978-0-521-45714-9。
\item 卢曼，J.G.;罗素，C。T.(1997)。“金星: 磁场和磁层”。在雪莉，J。H、; Fainbridge，R。W. (编辑)。行星科学百科全书。纽约:查普曼和霍尔。pp。 905-907。ISBN 978-1-4020-4520-2。已存档从原来的2010年7月14日。已检索7月19日2006。
\item 史蒂文森，D。J.(2003年3月15日)。“行星磁场”(PDF)。地球和行星科学快报。208(1-2):1-11.Bibcode:2003E & psl.208。... 1s。doi:10.1016/S0012-821X(02)01126-3。已存档(PDF)从原始的2017年8月16日。已检索11月6日2018。
\item 弗朗西斯·尼莫 (2002年11月)。“为什么金星没有磁场？”(PDF)。地质学。30(11):987-990。Bibcode:2002Geo .... 30 .. 987N。doi:10.1130/0091-7613(2002)030<0987:WDVLAM>2.0.CO;2。ISSN0091-7613。已存档(PDF)从2018年10月1日的原件。已检索6月28日2009年。  
\item Konopliv, A.美国；约德，C。F.(1996)。"金星人k2来自麦哲伦和PVO跟踪数据的潮汐爱情数字 “。地球物理研究快报。23(14):1857-1860年。Bibcode:1996年乔治。。23.1857k。doi:10.1029/96GL01589。
\item 雅各布森，塞思·A；鲁比，大卫·C；赫伦德，约翰; 莫比德利，亚历山德罗; 中岛，三木 (2017)。“地球和金星核心的形成、分层和混合”。地球和行星科学快报。474。Elsevier BV: 375。arXiv:1710.01770。Bibcode:2017E & PSL.474..375J。doi:10.1016/j.epsl.2017.06.023。S2CID119487513。 
\item Svedhem, Håkan; Titov, Dmitry V.;泰勒，弗雷德里克·W；Witasse，Olivier (2007年11月)。“金星是一个更像地球的行星”。自然。450(7170):629-632。Bibcode:2007 Natur.450..629S。doi:10.1038/nature06432。PMID18046393。S2CID1242297。  
\item O'Rourke，约瑟夫; Gillmann，塞德里克; Tackley，保罗 (2019年4月)。金星上的古代发电机和现代地壳残磁的前景。第21届EGU大会，EGU2019，会议记录于2019年4月7日至12日在奥地利维也纳举行。Bibcode:2019EGUGA..2118876O。18876。
\item 多纳休，T。M、；霍夫曼，J。H、; 霍奇斯，R。R、；华生，A.J.(1982)。“金星是湿的: 氘与氢之比的测量”。科学。216(4546):630-633。Bibcode:1982Sci...216 .. 630D。doi:10.1126/science.216.4546.630。ISSN0036-8075。PMID17783310。S2CID36740141。   
\item Lebonnois，塞巴斯蒂安; 舒伯特，杰拉尔德 (2017年6月26日)。“金星的深层大气以及密度驱动的CO2和N2分离的可能作用”(PDF)。自然地球科学。10(7)。施普林格科学与商业媒体有限责任公司:473-477。Bibcode:2017NatGe .. 10 .. 473L。doi:10.1038/ngeo2971。ISSN1752-0894。S2CID133864520。已存档(PDF)从2019年5月4日开始。已检索8月11日2023。   
\item 泰勒，弗雷德里克W.(2014)。“金星: 大气”。在蒂尔曼，斯波恩; 布鲁尔，多丽丝; 约翰逊，T。五、 (编辑)。太阳系百科全书。牛津:爱斯唯尔科学与技术。ISBN 978-0-12-415845-0。已存档从原件到2021年9月29日。已检索1月12日2016。
\item “维纳斯: 事实与数字”。NASA.存档自原来的2006年9月29日。已检索4月12日2007。
\item "维纳斯"。凯斯西储大学。2006年9月13日。存档自原来的2012年4月26日。已检索12月21日2011年。
\item 刘易斯，约翰S.(2004)。太阳系的物理和化学(第二版)。学术出版社。p。 463。ISBN 978-0-12-446744-6。
\item Newitz，Annalee (2013年12月11日)。“这是卡尔·萨根关于气候变化危险的原创文章”。Gizmodo。已存档从原件到2021年9月3日。已检索9月3日2021。
\item 杰弗里·A·兰迪斯.(2003)。“金星的殖民化”。AIP会议记录。第654卷。AIP. pp. 1193-1198。doi:10.1063/1.1541418。
\item Querol，Ricardo de (2023年4月28日)。"“宇宙”: 卡尔·萨根是对的 (我们没有太注意他)“。EL PAÍS English。已检索3月12日2025。
\item 哈利迪，亚历克斯·N.(2013年3月15日)。“地球行星中挥发物的起源”。Geochimica et Cosmochimica Acta。105:146-171。Bibcode:2013 GeCoA.105..146H。doi:10.1016/j.gca.2012.11.015。ISSN0016-7037。已存档从原件到2021年9月29日。已检索7月14日2020。 
\item Owen，Tobias; Bar-nun，Akiva; Kleinfeld，Idit (1992年7月)。“金星，地球和火星大气中重惰性气体的可能彗星起源”。自然。358(6381):43-46.Bibcode:1992年自然358. .. 43O。doi:10.1038/358043a0。ISSN1476-4687。PMID11536499。S2CID4357750。已存档从原件到2021年9月29日。已检索7月14日2020。   
\item 罗伯特·佩平.(1991年7月1日)。“关于地球大气和气象挥发物的起源和早期演化”。伊卡洛斯。92(1):2-79.Bibcode:1991年伊卡尔。92。2p。doi:10.1016/0019-1035(91)90036-S。ISSN0019-1035。 
\item Namiki，Noriyuki; 所罗门，肖恩C。(1998)。“氩和氦的火山脱气和金星上地壳生产的历史”。地球物理研究杂志: 行星。103(E2):3655-3677。Bibcode:1998JGR...103.3655N。doi:10.1029/97JE03032。ISSN2156-2202。 
\item 约瑟夫·G·奥罗克；Korenaga，Jun (2015年11月1日)。“金星的热演化与氩气脱气”。伊卡洛斯。260:128-140.Bibcode:2015年Icar .. 260 .. 128O。doi:10.1016/j.icarus.2015.07.009。ISSN0019-1035。 
\item 理查德·恩斯特 (2022年11月3日)。“金星曾经更像地球，但气候变化使它无法居住”。对话。已存档从2023年4月21日开始。已检索4月21日2023。
\item 路，M.J.;Del Genio，安东尼D.(2020)。“金星宜居气候情景: 通过时间和应用对金星进行建模，以缓慢旋转类似金星的系外行星”。地球物理研究杂志: 行星。125(5) e2019JE006276。美国地球物理联合会 (AGU)。arXiv:2003.05704。Bibcode:2020JGRE..12506276W。doi:10.1029/2019je006276。ISSN2169-9097。 
\item 路，M.J.;安东尼·D·德尔·吉尼奥；Kiang, Nancy Y.;琳达·索尔；Grinspoon，David H.; Aleinov，Igor; Kelley，Maxwell; Clune，Thomas (2016年8月28日)。金星是我们太阳系的第一个可居住的世界吗？。地球物理研究快报。43(16)。美国地球物理联合会 (AGU):8376-8383。arXiv:1608.00706。Bibcode:2016乔治。。43.8376W。doi:10.1002/2016gl069790。ISSN0094-8276。PMC5385710。PMID28408771。   
\item 格林斯彭，大卫·H。；布洛克，M。A.(2007年10月)。“在金星上寻找过去海洋的证据”。美国天文学会公报。39: 540。Bibcode:2007DPS....39.6109G。
\item Steigerwald，比尔 (2022年11月2日)。NASA研究: 大规模火山活动可能改变了古代金星的气候。NASA。已存档从2023年5月10日开始。已检索5月5日2023。
\item 卡斯特，J。F.(1988)。“失控和潮湿的温室大气以及地球和金星的演变”。伊卡洛斯。74(3):472-494。Bibcode:1988Icar...74 .. 472K。doi:10.1016/0019-1035(88)90116-9。PMID11538226。已存档从2019年12月7日的原件。已检索6月25日2019。 
\item 蒂尔曼，诺拉·泰勒 (2018年10月18日)。“金星的气氛: 组成，气候和天气”。Space.com。已存档从2023年5月9日开始。已检索5月9日2023。
\item Mullen，Leslie (2002年11月13日)。“金星云殖民地”。天体生物学杂志。存档自原来的2014年8月16日
\item 杰弗里·A·兰迪斯.(2003年7月)。“天体生物学: 金星的情况”(PDF)。英国星际学会杂志。56(7-8):250-254。Bibcode:2003JBIS...56..250L。NASA/TM-2003-212310。存档自原来的(PDF)2011年8月7日。
\item 科克尔，查尔斯S.(1999年12月)。“金星上的生命”。行星与空间科学。47(12):1487-1501.Bibcode:1999P & SS。47.1487c。doi:10.1016/S0032-0633(99)00036-7。
\item 帕特尔，M.R.；梅森，j.p.；诺德海姆，T.A.；Dartnell, L.R.(2022年)。“金星上潜在的空中生物圈的限制: II。紫外线辐射”(PDF)。伊卡洛斯。373114796。爱思唯尔BV。Bibcode:2022年Icar。。37314796p。doi:10.1016/j.icarus.2021.114796。ISSN0019-1035。S2CID244168415。   
\item 赫布斯特，康斯坦丁; 班纳克，萨沙; 阿特里，迪米特拉; 诺德海姆，汤姆A.(2019年12月24日)。“在可居住性的背景下重新审视宇宙射线引起的金星辐射剂量”。天文学与天体物理学。633。EDP科学: A15.arXiv:1911.12788。Bibcode:2020A & A...633A..15H。doi:10.1051/0004-6361/201936968。ISSN0004-6361。S2CID208513344。  
\item 纳迪亚·德雷克(2020年9月14日)。“金星上可能的生命迹象激起了激烈的辩论”。国家地理。存档自原来的2020年9月14日。已检索9月14日2020。
\item 格里夫斯，J。美国；理查兹，A。M。美国；贝恩斯，W.；Rimmer，P。B.;佐川，H; 克莱门茨，D。L.;西格，S.；Petkowski, J.J.;苏萨-席尔瓦，克拉拉; 兰詹，苏克里特; 德拉贝克-曼德，艾米丽; 弗雷泽，海伦·J。；卡特赖特，安娜贝尔; 穆勒-沃达格，英戈; 詹，朱昌; 弗里伯格，Per; 库尔森，伊恩; 李，埃丽莎; 霍格，吉姆 (2020)。“金星云层中的磷化氢气体”。自然天文学。5(7):655-664。arXiv:2009.06593。Bibcode:2021年纳塔斯。5。655克。doi:10.1038/s41550-020-1174-4。S2CID221655755。已存档从原件到2020年9月14日。已检索9月14日2020。 
\item 安德鲁·P·林考斯基；梅多斯，维多利亚S.；脆，大卫; 阿金斯，亚历克斯B；爱德华·W·施维特曼；阿尼，吉亚达N.；黄，迈克尔·L；保罗·G·斯特菲斯；Parenteau, M.尼基; 多马加尔-戈德曼，肖恩 (2021年)。“声称检测PH3在金星的云层中与中层SO2一致。天体物理学杂志。908(2): l44。arXiv:2101.09837。Bibcode:2021ApJ...908L..44L。doi:10.3847/2041-8213/abde47。S2CID231699227。 
\item 比尔，阿比盖尔 (2020年10月21日)。“对金星上潜在生命迹象的更多怀疑”。新科学家。doi:10.1016/S0262-4079(20)31910-2。S2CID229020261。已存档从原件到2021年12月26日。已检索1月29日2023。 
\item Snellen，我.A.G.;Guzman-Ramirez, L.;Hogerheijde, M.R、；Hygate, A.P。S、; 范德塔克，F。F.S。(2020年12月)。“金星267 GHz ALMA观测的重新分析”。天文学与天体物理学。644: l2。arXiv:2010.09761。Bibcode:2020A & A...644L...2S。doi:10.1051/0004-6361/202039717。S2CID224803085。已存档从2022年1月16日开始。已检索1月29日2023。 
\item 莫什金，B.E.;Ekonomov, A.P.;Golovin，Iu。M。(1979)。“金星表面的灰尘”。Kosmicheskie Issledovaniia (宇宙研究)。17(2):280-285。Bibcode:1979年科斯雷 .. 17 .. 232米。
\item “关于金星云的新细节揭示了”。欧洲航天局。2008。
\item “酸云和闪电”。欧洲航天局。
\item M.L.Delitsky和K.H.贝恩斯 (2015年11月)。“金星云中的化学: 含水液滴的硫酸反应和冻结行为”。AAS/行星科学部会议摘要 #47。47。美国天文学会。Bibcode:2015DPS....4721702D。
\item Krasnopolsky, V.A.;帕尔舍夫，V.A.(1981)。“金星大气的化学成分”。自然。292(5824):610-613。Bibcode:1981 Natur.292..610K。doi:10.1038/292610a0。S2CID4369293。 
\item 克拉斯诺波尔斯基，弗拉基米尔A.(2006)。“金星大气和云的化学成分: 一些尚未解决的问题”。行星与空间科学。54(13-14):1352-1359。Bibcode:2006P & SS。54.1352K。doi:10.1016/j.pss.2006.04.019。
\item 玛格丽特·戴维斯 (2021年7月14日)。“为什么金星如此明亮？这是它与地球的距离，高度反射的云层如何影响它。。科学时报。已存档从原始的2022年12月13日。已检索6月11日2023。
\item “金星和地球: 天壤之别-金星过境博客”。ESA博客导航器-活动ESA博客的导航器页面。2012年5月31日。已存档从2023年6月11日开始。已检索6月11日2023。
\item “金星的揭幕: 炎热和令人窒息”。科学新闻。109(25):388-389。1976年6月19日doi:10.2307/3960800。JSTOR3960800。每平方米100瓦。..14，000勒克斯...对应于...白天乌云密布 
\item Sánchez-lavega，奥古斯丁; Lebonnois，塞巴斯蒂安; Imamura，Takeshi; Read，Peter; Luz，David (2017年11月)。“金星的大气动力学”。空间科学评论。212(3-4):1541-1616。doi:10.1007/s11214-017-0389-x。ISSN0038-6308。 
\item Rossow, W.B、; 德尔·吉尼奥，A。D.;艾希勒，T.(1990)。“来自先锋金星OCPP图像的云追踪风”。大气科学杂志。47(17):2053-2084。Bibcode:1990JAtS...47.2053R。doi:10.1175/1520-0469(1990)047<2053:CTWFVO>2.0.CO;2。ISSN1520-0469。 
\item Antonio francés-monerris等人 (2022年7月30日)。“光化学和热化学路径 -- 通往S的途径2和金星大气中的多硫形成“。自然通讯。13(1) 4425。doi:10.1038/s41467-022-32170-x。PMC9338966。PMID35907911。  
\item Basilevsky, A.T.;头，J.W。(2003)。“金星的表面”。物理学进展报告。66(10):1699-1734.Bibcode:2003RPPh。66.1699b。doi:10.1088/0034-4885/66/10/R04。S2CID13338382。已存档从原件到2021年9月29日。已检索12月2日2019。 
\item 拉尔夫·D·洛伦兹；鲁尼，乔纳森一世；保罗·G·威瑟斯；麦凯，克里斯托弗·P。(2001年2月1日)。“泰坦，火星和地球: 纬向热传输产生的熵”(PDF)。地球物理研究快报。28(3)。艾姆斯研究中心亚利桑那大学月球与行星实验室:415-418。Bibcode:2001乔治。。28 .. 415L。doi:10.1029/2000GL012336。S2CID15670045。已存档(PDF)从2018年10月3日的原件。已检索8月21日2007。  
\item "星际季节"。NASA科学。NASA.2000年6月19日。已存档从原件到2021年4月14日。已检索4月14日2021。
\item 麦吉尔，G。E.;Stofan，E。R、；Smrekar, S.E.(2010)。“金星构造”。在Watters，T.R、；舒尔茨，R。A. (编辑)。行星构造。剑桥大学出版社。pp. 81-120。ISBN 978-0-521-76573-2。已存档从2016年6月23日的原件。已检索10月18日2015。
\item 奥滕，卡罗琳·琼斯 (2004)。"“重金属” 金星上的雪是硫化铅。圣路易斯华盛顿大学。已存档从2008年4月15日的原件。已检索8月21日2007。
\item Lee，Yeon Joo; Jessup，kandis-lea; Perez-hoyos，圣地亚哥; Titov，Dmitrij V.；Lebonnois, Sebastien; Peralta, Javier; Horinouchi, Takeshi; Imamura, Takeshi; Limaye, Sanjay; Marcq, Emmanuel; Takagi, Masahiro; Yamazaki, Atsushi; Yamada, Manabu; Watanabe, Shigeto; 村上，申亚; Ogohara，Kazunori; 麦克林托克，威廉M.；霍尔斯克劳，格雷戈里; 罗马，安东尼 (2019年8月26日)。“金星快车，明月，信使和哈勃太空望远镜观测到的金星365 nm反照率的长期变化”。天文杂志。158(3): 126。arXiv:1907.09683。Bibcode:2019AJ....158..126L。doi:10.3847/1538-3881/ab3120。S2CID198179774。
\item 拉尔夫·D·洛伦兹.(2018年6月20日)。“金星上的闪电探测: 批判性评论”。地球与行星科学进展。5(1) 34。Bibcode:2018PEPS....5...34L。doi:10.1186/s40645-018-0181-x。ISSN2197-4284。 
\item Kranopol'skii, V.A.(1980)。“根据卫星获得的信息，金星上的闪电维纳拉9和10”。宇宙研究。18(3):325-330。Bibcode:1980CosRe .. 18 .. 325K。
\item 罗素，C。T.;菲利普斯，J。L.(1990)。《灰烬之光》。空间研究进展。10(5):137-141.Bibcode:1990年AdSpR。。10 e.137R。doi:10.1016/0273-1177(90)90174-X。已存档从原来的2015年12月8日。已检索9月10日2015。
\item "维纳拉12下降艇”。国家空间科学数据中心。NASA.已存档从原件到2019年5月23日。已检索9月10日2015。
\item 罗素，C。T.;张，T。L.;德尔瓦，M.；马格尼斯，W.；Strangeway, R.J.;魏，h.y。(2007年11月)。“从电离层中的惠斯勒模式波推断出金星上的闪电”(PDF)。自然。450(7170):661-662。Bibcode:2007 Natur.450..661R。doi:10.1038/nature05930。PMID18046401。S2CID4418778。存档自原来的(PDF)2016年3月4日。已检索9月10日2015。   
\item 手，埃里克 (2007年11月)。“来自金星的欧洲任务报告”。自然(450):633-660。doi:10.1038/news.2007.297。S2CID129514118。 
\item 工作人员 (2007年11月28日)。“金星提供地球气候线索”。BBC新闻。已存档从2009年1月11日的原件。已检索11月29日2007。
\item “欧空局发现金星也有臭氧层”。欧洲航天局。2011年10月6日。已存档从2012年1月27日的原件。已检索12月25日2011年。
\item “当一颗行星表现得像彗星时”。欧洲航天局。2013年1月29日已存档从2019年5月2日的原件。已检索1月31日2013。
\item 克莱默，米里亚姆 (2013年1月30日)。“金星可以有 '彗星般的' 大气层”。Space.com。已存档从2019年5月3日的原件。已检索1月31日2013。
\item 福原哲也; Futaguchi，Masahiko; 桥本，乔治L.；Horinouchi, Takeshi; Imamura, Takeshi; Iwagaimi, Naomoto; Kouyama, Toru; Murakami, Shin-ya; Nakamura, Masato; Ogohara, Kazunori; Sato, Mitsuteru; Sato, Takao M.;铃木，诚; 田口，诚; 高木，精工; 上野，Munetaka; 渡边，茂东; 山田，Manabu; 山崎，Atsushi (2017年1月16日)。“金星大气中的大静止重力波”。自然地球科学。10(2):85-88。Bibcode:2017NatGe .. 10...85F。doi:10.1038/ngeo2873。
\item Rincon，Paul (2017年1月16日)。“金星波可能是太阳系最大的”。BBC新闻。已存档从2017年1月17日的原件。已检索1月17日2017。
\item Chang，Kenneth (2017年1月16日)。金星微笑着，神秘的波浪穿过它的大气层。纽约时报。已存档从2017年7月15日开始。已检索1月17日2017。
\item “Venera 14: 苏联科学院: 免费下载，借用和流式传输: 互联网档案馆”。互联网档案。2016年10月23日。已检索6月2日2025。
\item 凯恩，S.R、；Vervoort, P.;霍纳，J.；Pozuelos，P。J.(2020年9月)。“木星的迁移会加速金星的大气演化吗？”。行星科学杂志。1(2):42-51。arXiv:2008.04927。Bibcode:2020PSJ.....1...42K。doi:10.3847/PSJ/abae63。
\item “金星上一天的长度总是在变化-空间”。EarthSky。2021年5月5日。已存档从原始的2023年4月28日。已检索4月28日2023。
\item 让-吕克·玛戈特; 唐纳德·B·坎贝尔；乔恩·D·乔尔吉尼；Jao, Joseph S.;劳伦斯·G·斯内德克；弗兰克·D·吉戈；Bonsall，琥珀 (2021年4月29日)。“金星的自旋状态和转动惯量”。自然天文学。5(7):676-683。arXiv:2103.01504。Bibcode:2021年纳塔斯。5。676米。doi:10.1038/s41550-021-01339-7。ISSN2397-3366。 
\item “金星会换档吗？”。金星快车。欧洲航天局。2012年2月10日。已存档从原来的2016年1月24日。已检索1月7日2016。
\item “太空主题: 比较行星”。行星社会。存档自原来的2006年2月18日。已检索1月12日2016。
\item 塞尔日·布鲁尼耶(2002)。太阳系航行。邓洛普翻译，风暴。剑桥大学出版社。第40页。ISBN 978-0-521-80724-1。已存档从原件到2020年8月3日。已检索9月17日2017。
\item 科雷亚，亚历山大·C。M、；Laskar，Jacques; De Surgy，Olivier n é ron (2003年5月)。“金星自旋的长期演变，第一部分: 理论”(PDF)。伊卡洛斯。163(1):1-23.Bibcode:2003年Icar。。163 .... 1C。doi:10.1016/S0019-1035(03)00042-3。已存档(PDF)从原件到2019年9月27日。已检索9月9日2006。
\item Laskar，雅克; De Surgy，Olivier n é ron (2003)。“金星自旋的长期演变，第二部分: 数值模拟”(PDF)。伊卡洛斯。163(1):24-45。Bibcode:2003年Icar。。163。24C。doi:10.1016/S0019-1035(03)00043-5。已存档(PDF)从2019年5月2日的原件。已检索9月9日2006。
\item 黄金，T.；索特，S.(1969)。“大气潮汐与金星的共振旋转”。伊卡洛斯。11(3):356-66.Bibcode:1969年Icar。11 .. 356克。doi:10.1016/0019-1035(69)90068-2。
\item 夏皮罗，我.I.;坎贝尔，D。B.;De Campli, W.M。(1979年6月)。“金星的非共振旋转”。天体物理学杂志。230:L123-L126。Bibcode:1979ApJ...230L.123S。doi:10.1086/182975。
\item 斯科特·S·谢泼德；查德威克·特鲁希略。(2009年7月)。“金星卫星调查”。伊卡洛斯。202(1):12-16.arXiv:0906.2781。Bibcode:2009年Icar .. 202... 12s。doi:10.1016/j.icarus.2009.02.008。S2CID15252548。 
\item Mikkola, S.;布拉塞尔，R.；维格特，P.；Innanen，K。(2004年7月)。“小行星2002 VE68: 金星的准卫星”。皇家天文学会月刊。351(3): l63。Bibcode:2004MNRAS.351L..63M。doi:10.1111/j.1365-2966.2004.07994.x。
\item De la Fuente Marcos，Carlos; De la Fuente Marcos，ra ú l (2012年11月)。“关于2002 VE68的动态演变”。皇家天文学会月刊。427(1):728-39。arXiv:1208.4444。Bibcode:2012MNRAS.427..728D。doi:10.1111/j.1365-2966.2012.21936.x。S2CID118535095。 
\item De la Fuente Marcos，Carlos; De la Fuente Marcos，ra ú l (2012年11月)。“关于2002 VE68的动态演变”。皇家天文学会月刊。427(1):728-39。arXiv:1208.4444。Bibcode:2012MNRAS.427..728D。doi:10.1111/j.1365-2966.2012.21936.x。S2CID118535095。 
De la Fuente Marcos，卡洛斯; De la Fuente Marcos，劳尔 (2013年6月)。“小行星2012 XE133: 金星的短暂伴侣”。皇家天文学会月刊。432(2):886-93。arXiv:1303.3705。Bibcode:2013MNRAS.432..886D。doi:10.1093/mnras/stt454。S2CID118661720。 
\item 乔治·穆瑟 (2006年10月10日)。双重撞击可以解释为什么金星没有月亮。科学美国人。已存档从2007年9月26日的原件。已检索1月7日2016。
\item Tytell，David (2006年10月10日)。“为什么金星没有月亮？”。天空和望远镜。已存档从原来的2016年10月24日。已检索1月7日2016。
\item 莎拉·弗雷泽 (2021年4月16日)。NASA的帕克太阳探测器看到金星轨道尘埃环。NASA。已存档从2022年8月22日开始。已检索1月21日2023。
\item 加纳，罗布 (2019年3月12日)。“科学家在太阳系中筛选尘埃后发现了什么”。NASA。已存档从2023年1月1日开始。已检索1月21日2023。
\item 雷姆，杰里米 (2021年4月15日)。“帕克太阳探测器捕获了金星轨道尘埃环的第一个完整视图”。JHUAPL。已存档从2023年1月21日开始。已检索1月21日2023。
\item 巴兹索，A.；Eybl, V.;德沃夏克，R.；皮拉特-洛辛格，E.；Lhotka，C。(2010)。“金星和地球之间近平均运动共振的调查”。天体力学与动力天文学。107(1):63-76。arXiv:0911.2357。Bibcode:2010年CeMDA.107. .. 63B。doi:10.1007/s10569-010-9266-6。S2CID117795811。 
\item 奥特威尔，盖伊 (2022年1月7日)。“金星的5个花瓣及其8年周期”。EarthSky。已存档从2022年2月22日开始。已检索3月17日2022年。
\item 蒂姆·哈福德 (2019年1月11日)。“BBC广播4台-或多或少，糖，户外游戏和行星”。英国广播公司。已存档从原始的2019年1月12日。已检索10月30日2019。奥利弗·霍金斯，或多或少是校友和统计传奇，为我们写了一些代码，计算出过去50年里每一天哪颗行星离地球最近，然后将结果发送到大卫A.罗瑟里,开放大学行星地球科学教授。
\item “Solex 11预测的金星接近地球”。存档自原来的2012年8月9日。已检索3月19日2009年。生成的数字Solex
\item “金星不是地球最近的邻居。今日物理(3) 30593。AIP发布。2019年3月12日。Bibcode:2019PhT。。2019c0593。。doi:10.1063/pt.6.3.20190312a。ISSN1945-0699。S2CID241077611。  
\item 斯托克曼，汤姆; 门罗，加布里埃尔; 科德纳，塞缪尔 (2019)。“金星不是地球最近的邻居 | 计算和模拟证实，平均而言，水星是离地球最近的行星-以及太阳系中的所有其他行星”。今日物理。美国物理研究所。doi:10.1063/PT.6.3.20190312a。
\item 日元，C.-W。L.(1986年8月)。“通过金星和水星重力辅助的弹道水星轨道飞行器任务”。天体动力学1985:1293-1308。Bibcode:1986年asdy.conf.1293Y。
\item 弗朗西斯·M·斯特姆斯；切割，艾略特 (1966年5月)。“通过与金星的近距离接触对1970年水星任务的轨迹分析”。宇宙飞船和火箭杂志。3(5):624-631。doi:10.2514/3.28505。ISSN0022-4650。 
\item “帕克太阳探测器甚至在发射之前就改变了游戏-NASA”。2018年10月4日。已检索6月23日2025。
\item "IV-1"。nss.org。已检索6月22日2025。
\item 伊森伯格，诺姆·R；麦克纳特，拉尔夫·L；鲁尼翁，柯比·D；保罗·K·伯恩；麦克唐纳，亚历山大 (2021年3月1日)。“新人类太空时代的金星探索”。宇航学报。180:100-104.Bibcode:2021年AcAau.180..100I。doi:10.1016/j.actaastro.2020.12.020。ISSN0094-5765。 
\item Petropoulos, Anastassios E.;詹姆斯·M·朗格斯基；邦菲利奥，尤金·P。(2000)。“从金星、地球和火星通过重力辅助到达木星的轨迹”。宇宙飞船和火箭杂志。37(6)。美国航空航天学会 (AIAA):776-783。Bibcode:2000JSpRo。。37 .. 776P。doi:10.2514/2.3650。ISSN0022-4650。 
\item 泰勒，克里斯 (2020年7月9日)。“欢迎来到云城: 去金星的理由，而不是火星”。可混搭。已存档从2022年10月21日开始。已检索10月21日2022年。
\item "行星际低潮"。科学任务理事会。2000年5月3日。存档自原来的2023年6月4日。已检索6月25日2023。
\item 狄金森，特伦斯 (1998)。NightWatch: 观察宇宙的实用指南。纽约州布法罗: 萤火虫图书。第134页。ISBN 978-1-55209-302-3。已存档从原件到2021年9月29日。已检索1月12日2016。
\item “开始观察早晨天空中最明亮的金星”。地球和天空。2025年4月15日。已检索4月21日2025。
\item 弗兰德斯，托尼 (2011年2月25日)。“在光天化日之下见金星!”。天空和望远镜。存档自原来的2012年9月11日。已检索1月11日2016。
\item Espenak，Fred (1996)。“金星: 十二年行星星历，1995-2006”。NASA参考出版物1349。NASA/戈达德太空飞行中心。存档自原来的2000年8月17日。已检索6月20日2006。
\item "识别ufo"。夜空网络。太平洋天文学会。已存档从2021年4月10日开始。已检索4月10日2021。
\item 请参见本JPL地平线星历计算。
\item 薰衣草，杰玛 (2023年3月26日)。“你需要什么设备才能看到和拍摄行星”。Space.com。已检索6月5日2024。
\item 戈因斯，大卫·兰斯 (1995年10月18日)。“新月金星望远镜观测前的推论证据”。Goines.net。已存档从原始的2021年5月4日。已检索4月19日2017。
\item “在光天化日之下观看金星”。www.fourmilab.ch。已存档从原件到2021年11月15日。已检索7月17日2023。
\item Chris Chatfield (2010)。“用肉眼看太阳系”。自然现象画廊。已存档从原件到2015年6月13日。已检索4月19日2017。
\item Gaherty，Geoff (2012年3月26日)。今天在白天的天空中可见的金星: 如何看待它。Space.com。已存档从原件到2017年4月19日。已检索4月19日2017。
\item “2004年和2012年的金星凌日”。NASA。2004年6月8日。已存档从2023年5月2日开始。已检索5月2日2023。
\item 尼古拉斯·科勒斯特伦 (1998)。“霍罗克斯和英国天文学的曙光”。伦敦大学学院。已存档从2013年6月26日的原件。已检索5月11日2012。
\item Hornsby, T.(1771)。太阳的视差量，由对1769年6月3日金星凌日的观测推论而来。皇家学会的哲学交易。61:574-579。doi:10.1098/rstl.1771.0054。S2CID186212060。已存档从2019年5月9日的原始。已检索1月8日2008。 
\item 理查德·伍利 (1969)。“库克船长和1769年的金星凌日”。伦敦皇家学会的笔记和记录。24(1):19-32。doi:10.1098/rsnr.1969.0004。ISSN0035-9149。JSTOR530738。S2CID59314888。   
\item 艾伦·博伊尔 (2012年6月5日)。“金星过境: 最后一刻指南”。NBC新闻。存档自原来的2013年6月18日。已检索1月11日2016。
\item Espenak，Fred (2004)。“金星凌日，六个千年目录: 公元前2000年至公元4000年”。太阳的凌日。NASA.已存档从原件到2012年3月19日。已检索5月14日2009年。
\item 多宾斯，托马斯A.(2022年2月22日)。帕克太阳探测器捕捉到金星夜间令人惊讶的图像。天空和望远镜。已检索3月14日2025。
\item “帕克太阳探测器的夜间观测: 对灰光现实的影响-英国天文协会”。英国天文协会-自1890年以来支持业余天文学家。2024年5月21日。已检索3月14日2025。
\item 鲍姆，R。M。(2000)。“神秘的金星灰白之光: 概述”。英国天文协会杂志。110: 325。Bibcode:2000JBAA..110..325B。
\item Cooley, Jeffrey L.(2008)。“Inana和 š ukaletuda: 苏美尔星界神话”。卡斯卡。5:161-172。ISSN1971-8608。已存档从原始的2019年12月24日。已检索12月28日2017。 
\item 萨克斯，A。(1974)。“巴比伦观测天文学”。伦敦皇家学会的哲学交易。276(1257):43-50。Bibcode:1974RSPTA.276...43S。doi:10.1098/rsta.1974.0008。S2CID121539390。 
\item 罗素·霍布森 (2009)。公元前第一个千年文本的确切传播(PDF)(博士)。悉尼大学,希伯来语，圣经和犹太研究系。已存档(PDF)从2012年2月29日的原件。已检索12月26日2015。
\item Enn Kasak, Raul Veede.了解古代美索不达米亚的行星。民俗卷。16。Mare k õ iva & Andres Kuperjanov，Eds。ISSN 1406-0957
\item 海姆贝尔，W。(1982)。“近东金星神灵目录”。西罗-美索不达米亚研究。4(3)。Undena出版物:9-22.
\item 约瑟夫·李约瑟(1959) [1959]。数学与天地科学。科学与文明在中国卷3。剑桥: 剑桥大学出版社。第398页。Bibcode:1959scc3。书 ..... N。ISBN 978-0-521-05801-8。
\item 老普林尼 (1991)。自然历史II:36-37。翻译: Healy，John F.Harmondsworth，米德尔塞克斯，英国:企鹅。pp。 15-16.
\item 沃尔特·伯克特 (1972)。古代毕达哥拉斯主义的知识和科学。哈佛大学出版社。第307页。ISBN 978-0-674-53918-1。已存档从2016年6月9日的原件。已检索12月28日2015。
\item 多宾，罗伯特 (2002)。“具有讽刺意味的是” Aeneid “1.374”。Mnemosyne。第四系列。55(6)。布里尔:736-738。doi:10.1163/156852502320880285。JSTOR4433390。 
\item 伯纳德·R·戈尔茨坦.(1972年3月)。“中世纪天文学中的理论和观察”。Isis。63(1): 39-47 [44]。Bibcode:1972Isis...63...39G。doi:10.1086/350839。S2CID120700705。 
\item “阿维森纳viii。数学和物理科学“。伊朗百科全书。已存档从原件到2020年2月20日。已检索3月4日2016。
\item 安萨里，S.M。Razaullah (2002年)。东方天文学史。1997年8月25日至26日在京都举行的第41委员会 (天文学史) 组织的国际天文学联合会第23届大会的联合讨论记录。天体物理学和空间科学图书馆。施普林格科学 + 商业媒体。第137页。ISBN 978-1-4020-0657-9。
\item Vaquero, J.M、；Vázquez, M.(2009)。通过历史记录的太阳。施普林格科学与商业媒体。第75页。ISBN 978-0-387-92790-9。已存档从原来的2016年11月26日。已检索5月18日2016。
\item 肯纳德，弗雷德里克 (2015)。思想实验: 哲学，物理学，伦理学，计算机科学和数学中的流行思想实验。Lulu.com。第113页。ISBN 978-1-329-00342-2。已存档从原来的2016年11月25日。已检索5月18日2016。
\item 保罗·帕尔米耶里 (2001年)。“伽利略与金星相位的发现”。天文学史杂志。21(2):109-129。Bibcode:2001JHA....32..109P。doi:10.1177/002182860103200202。S2CID117985979。 
\item 小费格利，B。(2003)。"维纳斯"在荷兰，海因里希·D；卡尔·K·图雷基安 (编辑)。地球化学论文。爱思唯尔.pp. 487-507.ISBN 978-0-08-043751-4。
\item 尼古拉斯·科勒斯特伦 (2004)。“威廉·克拉布特里的金星凌日观测”(PDF)。国际天文学联合会第196号学术讨论会。2004:34-40.Bibcode:2005tvnv.conf...34K。doi:10.1017/S1743921305001249。S2CID162838538。已存档(PDF)从2016年5月19日的原件。已检索5月10日2012。  
\item 马洛夫，米哈伊尔亚。(2004)。库尔茨，D。W. (编辑)。米哈伊尔·罗莫诺索夫和1761年过境期间金星大气的发现。金星凌日: 太阳系和银河系的新观点，IAU学术讨论会论文集 #196，2004年6月7日至11日在英国普雷斯顿举行。2004年卷。剑桥大学出版社。pp. 209-219。Bibcode:2005tvnv.conf .. 209米。doi:10.1017/S1743921305001390。
\item "米哈伊尔·瓦西里耶维奇·罗莫诺索夫"。大英百科全书在线。已存档从2008年7月25日开始。已检索7月12日2009年。
\item 罗素，H. N.(1899)。“金星的大气层”。天体物理学杂志。9:284-299。Bibcode:1899ApJ ..... 9 .. 284R。doi:10.1086/140593。S2CID123671250。 
\item 赫西，T。(1832)。“关于金星的旋转”。皇家天文学会月刊。2(11):78-126。Bibcode:1832MNRAS...2...78H。doi:10.1093/mnras/2.11.78d。已存档从原件到2020年7月11日。已检索8月25日2019。
\item 罗斯，F.E.(1928)。“金星的照片”。天体物理学杂志。68: 57。Bibcode:1928年4月。68...57R。doi:10.1086/143130。
\item 马茨，埃德温·P。小 (1934)。“维纳斯和生命”。热门天文学。42: 165。Bibcode:1934帕... 42.165m。
\item Mitchell，Don (2003)。“发明行星际探测器”。苏联对金星的探索。已存档从原件到2018年10月12日。已检索12月27日2007。
\item Mayer, C.H、; 麦卡洛，T。P.;Sloanaker, R.M。(1958年1月)。“金星在3.15厘米波长处的观测”。天体物理学杂志。127: 1。Bibcode:1958年apj。127。1米。doi:10.1086/146433。
\item 喷气推进实验室 (1962)。水手-维纳斯1962年最终项目报告(PDF)(报告)。SP-59。NASA.已存档(PDF)从原件到2014年2月11日。已检索7月7日2017。
\item 戈德斯坦，R。M、；卡彭特，R。L.(1963)。“金星的旋转: 根据雷达测量估算的周期”。科学。139(3558):910-911.Bibcode:1963年科学… 139 .. 910克。doi:10.1126/science.139.3558.910。PMID17743054。S2CID21133097。  
\item Mitchell，Don (2003)。“维纳斯的大气层”。苏联对金星的探索。已存档从原件到2018年9月30日。已检索12月27日2007。
\item 哈维，布莱恩 (2007)。俄罗斯行星探测历史、发展、遗产与展望。施普林格-实践。pp. 115-118.ISBN 978-0-387-46343-8。
\item “关于空间研委会第七工作组活动的报告”。空间研委会第十二次全体会议和第十届国际空间科学专题讨论会初步报告。捷克斯洛伐克布拉格:国家科学院。1969年5月11日至24日。第94页。
\item “科学: 从金星向前”。时间。1971年2月8日存档自原来的2008年12月21日。已检索1月2日2013。
\item 坎贝尔，D。B.;戴斯，R。B.;佩滕吉尔，G。H. (1976)。“金星的新雷达图像”。科学。193(4258):1123-1124。Bibcode:1976Sci。193.1123C。doi:10.1126/science.193.4258.1123。PMID17792750。S2CID32590584。  
\item 科林，L.；霍尔，C。(1977)。“先锋金星计划”。空间科学评论。20(3):283-306。Bibcode:1977SSRv...20 .. 283C。doi:10.1007/BF02186467。S2CID122107496。 
\item 威廉姆斯，大卫R.(2005年1月6日)。"先锋金星项目信息"。NASA/戈达德太空飞行中心。已存档从原始的2019年5月15日。已检索7月19日2009年。
\item 罗纳德·格里利巴特森，雷蒙德·M。(2007)。行星映射。剑桥大学出版社。第47页。ISBN 978-0-521-03373-2。已存档从原件到2021年9月29日。已检索7月19日2009年。
\item “欢迎来到伽利略轨道飞行器档案页面”。PDS大气节点。1989年10月18日。已存档从2023年4月11日开始。已检索4月11日2023。
\item 伊丽莎白·豪厄尔 (2014年12月16日)。“金星快车耗尽了气体; 任务结束，死亡观察飞船”。今天的宇宙。已存档从原件到2021年4月22日。已检索4月22日2021。
\item 哈特菲尔德，英里 (2022年2月9日)。帕克太阳探测器捕获金星表面的可见光图像。NASA。已存档从2022年4月14日开始。已检索4月29日2022年。
\item 伍德，B.E.;赫斯，P.；Lustig-Yaeger, J.;加拉格尔，B.；科万，D.；里奇，N.；斯滕堡，G.；Thernisien, A.;Qadri, S.N.;圣地亚哥，F.；佩拉尔塔，J.；阿尼，G.N.;Izenberg, N.R、；Vourlidas, A.;林顿，M.G.;霍华德，R。A.;Raouafi, N.E.(2022年2月9日)。“帕克太阳探测器对金星夜侧的成像”。地球物理研究快报。49(3) e2021GL096302。Bibcode:2022乔治。。4996302W。doi:10.1029/2021GL096302。PMC9286398。PMID35864851。  
\item 约瑟夫·G·奥罗克；科林·F·威尔逊；博雷利，麦迪逊E.；保罗·K·伯恩；杜穆林，卡罗琳; 盖尔，理查德; 居尔彻，安娜·J。P.;雅各布森，塞思·A；科拉布列夫，奥列格; 斯波恩，蒂尔曼; 路，M。J.;马特·韦勒; 弗朗西斯·韦斯托尔 (2023)。“金星，行星: 地球姊妹行星进化简介”。空间科学评论。219(1)。施普林格科学与商业媒体有限责任公司: 10。Bibcode:2023年SSRv。。219。10O。doi:10.1007/s11214-023-00956-0。高密度脂蛋白:20.500.11850/598198。ISSN0038-6308。S2CID256599851。  
\item Rao，Rahul (2020年7月7日)。科学家说，前往火星的宇航员应该先经过金星。Space.com。已存档从2023年4月24日开始。已检索4月24日2023。
\item 伊森伯格，诺姆·R；麦克纳特，拉尔夫·L；鲁尼翁，柯比·D；保罗·K·伯恩；麦克唐纳，亚历山大 (2021)。“新人类太空时代的金星探索”。宇航学报。180。爱思唯尔BV:100-104.Bibcode:2021年AcAau.180..100I。doi:10.1016/j.actaastro.2020.12.020。ISSN0094-5765。S2CID219558707。  
\item 阿雷东多，阿尼西亚; 霍奇斯，阿莫雷; 亚伯拉罕，雅各布·N。H.; 贝德福德，坎迪斯C.；本杰明·D·博特赖特；布兹，詹妮弗; 坎特罗尔，克莱顿; 克拉克，乔安娜; 欧文，安德鲁; 克里希纳穆尔蒂，悉达多; 玛加尼亚，利泽斯; 麦凯布，瑞安M；麦金托什，E。嘉莉; 诺维耶洛，杰西卡·L；佩莱格里诺，玛丽埃尔; 雷，克里斯汀; 斯蒂辛斯基，马歇尔J；彼得·韦格尔 (2022年7月1日)。“情人节: 一种新的前沿级长时间原位基于气球的航空机器人对金星的任务的概念”。行星科学杂志。3(7): 152。Bibcode:2022PSJ.....3..152A。doi:10.3847/PSJ/ac7324。ISSN2632-3338。 
\item 杰弗里·A·兰迪斯.(2003)。“金星的殖民化”。AIP会议记录。第654卷。pp。 1193-1198。doi:10.1063/1.1541418。存档自原来的2012年7月11日
\item "Архив фантастики"。Архив фантастики(俄语)。已存档从原始的2021年9月2日。已检索9月2日2021。
\item 巴德斯库，维奥雷尔; 扎克尼，克里斯，编辑。(2015)。内太阳系。斯普林格国际出版社。doi:10.1007/978-3-319-19569-8。ISBN 978-3-319-19568-1。
\item Tickle，格伦 (2015年3月5日)。“看看人类是否应该尝试殖民金星而不是火星”。笑鱿鱼。已存档从原件到2021年9月1日。已检索9月1日2021。
\item Warmflash，大卫 (2017年3月14日)。“金星云的殖民化: '表面主义' 是否模糊了我们的判断？”。视觉学习。已存档从2019年12月11日的原件。已检索9月20日2019。
\item 克拉克，斯图尔特 (2003年9月26日)。“金星的酸性云可能孕育生命”。新科学家。已存档从原来的2015年5月16日。已检索12月30日2015。
\item 雷德芬，马丁 (2004年5月25日)。“金星云” 可能孕育生命"。BBC新闻。已存档从原件到2020年9月16日。已检索12月30日2015。
\item 刘易斯·R·达特内尔；汤姆·安德烈·诺德海姆; 曼尼什·R·帕特尔；乔纳森·P·梅森；安德鲁·J·科茨；Jones，Geraint H. (2015年9月)。“金星上潜在的航空生物圈的限制: I.“宇宙射线”。伊卡洛斯。257:396-405.Bibcode:2015Icar .. 257 .. 396D。doi:10.1016/j.icarus.2015.05.006。
\item 卡尔·萨根;莫罗维茨，哈罗德·J.(1967年9月16日)。“金星云层中的生命？”。自然。215(5107):1259-1260。doi:10.1038/2161198a0。S2CID11784372。已存档从原件到2020年9月17日。已检索9月17日2020。 
\item 保罗·安德森 (2019年9月3日)。“微生物会影响金星的气候吗？”。地球和天空。已存档从2019年9月3日的原件。已检索9月3日2019。
\item 威廉·贝恩斯; 雅努斯·J·佩特科夫斯基；萨拉·西格; 苏克里特·兰扬; 克拉拉·索萨·席尔瓦; 保罗·B·里默；詹，朱昌; 格里夫斯，简·S；理查兹，安妮塔·M。S。(2021)。“金星上的磷化氢不能用常规过程来解释”。天体生物学。21(10):1277-1304。arXiv:2009.06499。Bibcode:2021年阿斯比奥。。21.1277b。doi:10.1089/ast.2020.2352。PMID34283644。S2CID221655331。  
\item 珀金斯，Sid (2020年9月14日)。“好奇和无法解释”。科学。已存档从原件到2020年9月14日。已检索9月14日2020。
\item 西格，萨拉; 佩特科夫斯基，雅努斯；高，彼得; 贝恩斯，威廉; 布莱恩，诺埃尔·C；兰詹，苏克里特; 格里夫斯，简 (2020年9月14日)。“金星低层大气雾霾作为干燥微生物生命的仓库: 金星航空生物圈持久性的拟议生命周期”。天体生物学。21(10):1206-1223.arXiv:2009.06474。doi:10.1089/ast.2020.2244。PMID32787733。S2CID221127006。  
\item 样本，伊恩 (2020年9月14日)。科学家在金星大气中发现与生命有关的气体。守护者。已存档从原件到2021年2月5日。已检索9月16日2020。
\item 库瑟，阿曼达 (2020年9月14日)。“美国宇航局局长呼吁在意外发现外星生命的暗示后优先考虑金星”。CNet。已存档从原件到2020年9月15日。已检索9月14日2020。
\item @ JimBridenstine (2020年9月14日)。“金星上的生命？”(鸣叫)-通过Twitter。
\item 菲尔普莱特 (2020年10月26日)。“更新: 地狱之上的生活？对金星磷化氢发现的严重怀疑“。Syfy.com。Syfy。已存档从原件到2020年10月29日。已检索10月26日2020。
\item Snellen，我.A.G.;Guzman-Ramirez, L.;Hogerheijde, M.R、；Hygate, A.P。S、; 范德塔克，F。F.S。(2020)。“金星267 GHz ALMA观测的重新分析”。天文学与天体物理学。644: l2。arXiv:2010.09761。Bibcode:2020A & A...644L...2S。doi:10.1051/0004-6361/202039717。S2CID224803085。 
\item 汤普森，M。A.(2021)。“金星的267 GHZ JCMT观测的统计可靠性: 没有明显的磷化氢吸收证据”。皇家天文学会每月通知: 信件。501(1):L18-L22。arXiv:2010.15188。Bibcode:2021MNRAS.501L..18T。doi:10.1093/mnrasl/slaa187。
\item 维拉纽瓦，杰罗尼莫; 科迪纳，马丁; 欧文，帕特里克; 德佩特，伊姆克; 巴特勒，布莱恩; 古威尔，马克; 米拉姆，斯蒂芬妮; 尼克松，康纳; 卢什茨库克，斯塔西亚；威尔逊，科林; 科夫曼，文森特; 刘齐，朱利亚诺; 法吉，萨拉; 福切斯，托马斯; 里皮，曼努埃拉; 科森蒂诺，理查德; 塞伦，亚历山大; 穆勒，阿里埃勒; 哈托格，保罗; 莫尔特，爱德华; 查恩利，史蒂夫; 阿尼，贾达; 曼德尔，阿维; 比弗，尼古拉斯; 凡达，安; 德·克莱尔，凯瑟琳; 科帕拉普，拉维 (2021年)。“独立分析没有证据表明金星大气中存在磷化氢”。自然天文学。5(7):631-635。arXiv:2010.14305。Bibcode:2021年纳塔斯。5。631V。doi:10.1038/s41550-021-01422-z。S2CID236090264。 
\item "火箭实验室探针"。金星云生活-麻省理工学院。2023年3月7日。已存档从2024年2月8日开始。已检索5月13日2023。
\item 国家研究委员会 (2006年)。评估金星任务的行星保护要求: 信函报告。国家科学院出版社。doi:10.17226/11584。ISBN 978-0-309-10150-9。已存档从原件到2015年7月17日。已检索1月19日2021。
\item 佐尔扎诺，玛丽亚·帕斯; 奥尔森·弗朗西斯，卡伦; 多兰，彼得·T。；Rettberg，Petra; Coustenis，雅典娜; Ilyin，Vyacheslav; Raulin，Francois; Shehhi，Omar Al; Groen，Frank; Grasset，Olivier; Nakamura，Akiko; Ballesteros，Olga Prieto; Sinibaldi，西尔维奥; 铃木，Yohey; 库马尔，Praveen; Kminek，Gerhard; 赫德曼，尼克拉斯; 藤本，Masaki; 扎伊采夫，马克西姆; 海斯，亚历克斯; 彭，静; 安曼尼托，埃莱奥诺拉; 穆斯汀，克里斯蒂安; 徐，干扬 (2023)。“空间研委会对金星太空飞行任务的行星保护要求”。空间研究中的生命科学。37:18-24.doi:10.1016/j.lssr.2023.02.001。
\item 杰弗里·A·兰迪斯.(2020年11月16日)。定居金星: 云中的城市？。美国航空航天学会。doi:10.2514/6.2020-4152。ISBN 978-1-62410-608-8。已检索9月11日2025。
\item 布莱克，杰里米; 绿色，安东尼 (1992)。古代美索不达米亚的神，恶魔和符号: 插图词典。大英博物馆出版社。pp. 108-109。ISBN 978-0-7141-1705-8。已存档从原件到2020年11月20日。已检索8月23日2020。
\item Nemet-nejat，Karen Rhea (1998)。古代美索不达米亚的日常生活。格林伍德。第203页。ISBN 978-0-313-29497-6。已检索2月22023。
\item Cooley, Jeffrey L.(2008)。“Inana和 š ukaletuda: 苏美尔星界神话”。卡斯卡。5:163-164.ISSN1971-8608。已存档从原始的2019年12月24日。已检索12月28日2017。 
\item 帕克，R。A.(1974)。“古埃及天文学”。伦敦皇家学会的哲学交易。A辑，数学和物理科学。276(1257)。皇家学会:51-65.Bibcode:1974RSPTA.276...51P。doi:10.1098/rsta.1974.0009。ISSN0080-4614。JSTOR74274。S2CID120565237。   
\item 庸医，约阿希姆·弗里德里希 (2019年5月23日)。“古埃及的行星”。牛津研究行星科学百科全书。牛津大学出版社。doi:10.1093/因此/9780190647926.013.61。ISBN 978-0-19-064792-6。
\item 卡特莫尔，彼得·约翰; 摩尔，帕特里克 (1997)。金星地图集。剑桥大学出版社。第9页。ISBN 978-0-521-49652-0。
\item "路西法"。大英百科全书在线。2020年1月24日存档自原来的2020年1月24日。已检索2月3日2023。
\item 马库斯·图利乌斯·西塞罗 (2005年9月12日)。De Natura Deorum。存档自原来的2005年9月12日。已检索2月3日2023。。
\item 阿特斯马，亚伦·J。"Eospheros & Hespheros"。Theoi.com。已存档从原始的2019年7月14日。已检索1月15日2016。
\item 达瓦·索贝尔(2005)。行星。哈珀出版。pp。 53-70。ISBN 978-0-14-200116-5。
\item 巴拉，Prem P.(2006)。印度教仪式，仪式，习俗和传统: 从A到Z论印度教的生活方式。Pustak Mahal。第29页。ISBN 978-81-223-0902-7。
\item 贝哈里，贝宾; 弗劳利，大卫 (2003)。吠陀占星术的神话与符号(第二版)。莲花出版社。pp. 65-74.ISBN 978-0-940985-51-3。
\item “阿芙罗狄蒂和爱的众神: 罗马维纳斯 (盖蒂别墅展览)”。盖蒂。已存档从2023年4月12日开始。已检索4月15日2023。
\item Jan Jakob Maria De Groot (1912)。中国的宗教: 大学主义。道家与儒学研究的关键。美国关于宗教史的讲座。第10卷。G.P。普特南的儿子。第300页。已存档从2011年7月22日开始。已检索1月8日2010。
\item 克鲁普，托马斯 (1992)。日本的数字游戏: 现代日本对数字的使用和理解。Routledge. pp. 39-40.ISBN 978-0-415-05609-0。
\item Hulbert，Homer Bezaleel (1909)。韩国的通过。双日，佩奇和公司。p。 426。已检索1月8日2010。
\item "Sao Kim - VOER"。越南开放教育资源。已存档从原来的2022年12月26日。已检索12月26日2022年。
\item Chumayel的书: Yucatec Maya的法律顾问书，1539-1638年。理查德·卢克斯顿.1995年 [1899]。第6页，第194页。ISBN 978-0-89412-244-6。
\item 苏珊·米尔布拉斯 (1999)。玛雅人的星神: 艺术，民间传说和日历中的天文学。德克萨斯州奥斯汀: 德克萨斯大学出版社。pp. 200-204，383。ISBN 978-0-292-79793-2。
\item 查尔斯·惠特尼.(1986年9月)。《文森特·梵高的天空》艺术史。9(3): 356。doi:10.1111/j.1467-8365.1986.tb00206.x。
\item 艾伯特·博伊姆(1984年12月)。梵高的星夜“物质的历史和历史的问题”(PDF)。艺术杂志: 88。已存档(PDF)从原件到2018年11月23日。已检索7月28日2018。
\item 米勒，罗恩 (2003)。维纳斯。《二十一世纪的书》，第12页。ISBN 978-0-7613-2359-4。
\item 迪克，史蒂文 (2001)。其他世界上的生命: 20世纪外星生命辩论。剑桥大学出版社。第43页。ISBN 978-0-521-79912-6。
\item 种子，大卫 (2005)。科幻小说的伴侣。布莱克韦尔出版。pp. 134-135.ISBN 978-1-4051-1218-5。已检索2月3日2023。
\item 肖特，G。D.(2005年12月22日)。“古代和现代的性符号: 它们的起源和谱系上的肖像学”。BMJ。331(7531):1509-1510。doi:10.1136/bmj.331.7531.1509。ISSN0959-8138。PMC1322246。PMID16373733。   
\item 斯蒂恩，威廉T.(1961年8月17日)。“生物学的男性和女性符号”。新科学家(248):412-413。存档自原来的2021年11月23日
\item 斯蒂恩，威廉T.(1968年5月)。“生物学中男性和女性符号的起源”。分类群。11(4):109-113。Bibcode:1962年Taxon .. 11 .. 109S。doi:10.2307/1217734。JSTOR1217734。S2CID87030547。  
\item Brammer，John Paul (2020年2月10日)。“爱/恨读: '男人来自火星，女人来自金星，' 重访”。副。已存档从2023年4月17日开始。已检索4月17日2023。
\item 卡尔·G·利格曼。(2004)。符号: 西方符号和表意符号百科全书。HME出版。第228页。ISBN 978-91-972705-0-2。
\end{enumerate}
\subsection{延伸阅读}
\begin{itemize}
\item O'Rourke, Joseph G.; Wilson, Colin F.; Borrelli, Madison E.; Byrne, Paul K.; Dumoulin, Caroline; Ghail, Richard; Gülcher, Anna J. P.; Jacobson, Seth A.; Korablev, Oleg; Spohn, Tilman; Way, M. J.; Weller, Matt; Westall, Frances（2023 年 2 月 6 日）。《金星，这颗行星：地球姊妹行星演化导论》。《空间科学评论》（Space Science Reviews）。219（1）：10。Bibcode:2023SSRv..219...10O。doi:10.1007/s11214-023-00956-0。hdl:20.500.11850/598198。ISSN 1572-9672。2025 年 2 月 6 日检索。
\item Widemann, Thomas；等（2023 年 10 月）。《金星随时间演化：关键科学问题、选定的任务设想与未来研究》。《空间科学评论》（Space Science Reviews）。219（7）：56。Bibcode:2023SSRv..219...56W。doi:10.1007/s11214-023-00992-w。hdl:10852/109541。
\end{itemize}
\subsection{外部链接}
\begin{itemize}
\item Venus profile at NASA's Solar System Exploration site
NASA 太阳系探索网站的金星概览页面
\item Missions to Venus Archived 10 February 2015 at the Wayback Machine and Image catalogue Archived 29 March 2015 at the Wayback Machine at the National Space Science Data Center
国家空间科学数据中心的“金星任务”页面（2015 年 2 月 10 日存档）及图像目录（2015 年 3 月 29 日存档）
\item Soviet Exploration of Venus and Image catalogue at Mentallandscape.com Mentallandscape.com 上的“苏联金星探测”与图像目录
\item Image catalogue from the Venera missions “金星（Venera）计划”图像目录

\item Venus page at The Nine Planets The Nine Planets 网站的金星页面

\item Transits of Venus Archived 19 March 2012 at the Wayback Machine at NASA.gov NASA.gov 上的“金星凌日”页面（2012 年 3 月 19 日存档）
\item Geody Venus, a search engine for surface features Geody Venus：金星表面地形搜索引擎
\item Interactive 3D gravity simulation of the pentagram that the orbit of Venus traces when Earth is held fixed at the centre of the coordinate system交互式 3D 重力模拟：在将地球固定于坐标系中心时，展示金星轨道绘出的五角形图案
\end{itemize}
\subsubsection{制图资源}
\begin{itemize}
\item Map-a-Planet: Venus by the U.S. Geological Survey
美国地质调查局（USGS）的“Map-a-Planet：金星”
\item Gazetteer of Planetary Nomenclature: Venus by the International Astronomical Union 国际天文学联合会的《行星命名公报：金星》
\item Venus crater database by the Lunar and Planetary Institute 月球与行星研究所提供的金星撞击坑数据库
\item Map of Venus Archived 6 November 2015 at the Wayback Machine by Eötvös Loránd University厄特沃什·罗兰大学的金星地图（2015 年 11 月 6 日存档）
\item Google Venus 3D, interactive map of the planet Google Venus 3D：金星互动式三维地图
\end{itemize}

